\documentclass[12pt]{article}

% ------------------------------------------------------------
% Essential Packages
% ------------------------------------------------------------
\usepackage[T1]{fontenc}
\usepackage[utf8]{inputenc}
\usepackage{microtype}

\usepackage{amsmath, amssymb, amsthm, mathtools}
\usepackage{geometry}
\geometry{margin=1in}

\usepackage{bm}
\usepackage[title]{appendix}

% ------------------------------------------------------------
% Custom Theorem / Statement Environment
% ------------------------------------------------------------
\newtheoremstyle{statementstyle}
  {6pt}{6pt}%
  {\itshape}{}%
  {\bfseries}{.}%
  {0.8em}{}%

\theoremstyle{statementstyle}
\newtheorem{statement}{Statement}[section]

% ------------------------------------------------------------
% Title Information
% ------------------------------------------------------------
\title{An Axiomatic Derivation of Instantaneous Rolling Forward Rate Dynamics via the Kolmogorov Forward Equation}

\author{Sachin Oza}
\date{March 2026}

% ------------------------------------------------------------
% Document Begins
% ------------------------------------------------------------
\begin{document}

\maketitle
\vspace{-1.0em}

% ------------------------------------------------------------
% Abstract
% ------------------------------------------------------------
\begin{abstract}
This paper develops an axiomatic framework for the dynamics of the instantaneous rolling forward rate (IRFR) and, by extension, the term structure of interest rates. Rather than postulating a stochastic process \emph{a priori}, it derives the admissible dynamics from four structural axioms encoding rolling-maturity invariance, equilibrium moment transport, and positivity at zero maturity. Under these restrictions, the IRFR is shown to follow a Pearson-type diffusion with affine drift and quadratic volatility, so that the diffusion structure arises endogenously from the axioms rather than being imposed exogenously.

The framework yields several concrete implications. The equilibrium forward curve is exponential, the zero-maturity equilibrium forward rate equals one half of the asymptotic equilibrium level, and the entire forward curve can be reconstructed from a single observed forward rate together with structural parameters. Solving the associated Kolmogorov Forward Equation further shows that the maturity-indexed equilibrium density is a shifted inverse-gamma law of order three, implying strict positivity and heavy tails. The model also delivers a constant market signal-to-noise ratio and a forward-curve transition structure that is analogous to covariance kernels in classical affine term-structure models. Taken together, these results provide a structural account of IRFR dynamics linking equilibrium behaviour, forward-curve evolution, and probabilistic structure.
\end{abstract}

% ------------------------------------------------------------
% 1. Introduction
% ------------------------------------------------------------

% --- Equilibrium terminology clarification ---
\paragraph{Equilibrium terminology.}
Throughout, we work with the maturity-indexed equilibrium regime. All distributional results are expressed as maturity-indexed equilibrium densities parameterized by $\tau = T-t$, rather than as limits in calendar time.

\section{Introduction}
\label{sec:introduction}

A central modelling choice in any term-structure framework is the 
selection of a rate-like quantity whose dynamics are both analytically 
tractable and economically meaningful. This paper adopts as its 
primitive the \emph{instantaneous rolling forward rate} (IRFR). 

Consider a forward-starting infinitesimal borrowing--lending contract 
initiated at time $t_{0}$ with zero initial value. At maturity 
$T - t_{0}$, the buyer pays an instantaneous interest rate 
$x_{t_{0},T}$ on a notional principal. This value defines the IRFR. 
Because the contract is continuously marked to market, its value 
process is, by construction,
\[
V(t) = x_{t,T},
\]
and over an infinitesimal interval $\Delta t$,
\[
\Delta V = x_{t+\Delta t,T} - x_{t,T}
= \Delta x_{t,T}.
\]

As time advances from $t$ to $t + \delta t$ while the maturity 
date $T$ remains fixed, the remaining maturity decreases from $T - t$ 
to $T - (t + \delta t)$. Consequently, the IRFR ``rolls'' from 
$x_{t,T}$ to $x_{t+\delta t,T}$. This rolling structure is a defining 
feature of the IRFR and distinguishes it from the classical 
Heath--Jarrow--Morton (HJM) forward rate, which is indexed by a fixed 
maturity date rather than remaining maturity.

In continuous time, the IRFR satisfies the fundamental bond-pricing 
relations
\[
\int_{t_{0}}^{T} x_{t_{0},S}\,dS = -\ln P(t_{0},T),
\qquad
\int_{t}^{T} x_{t,S}\,dS = -\ln P(t,T),
\]
where $P(t,T)$ denotes the zero-coupon bond price and $T-t$ is the 
current time-to-maturity.

The objective of this paper is to characterise the class of stochastic
processes governing the evolution of $x_{t,T}$ under four economically
and probabilistically natural axioms. These
axioms impose coherence across maturities and time, encode equilibrium
behaviour, and enforce positivity at the short end. Their consequences
are as follows:

\begin{itemize}
\item The drift and diffusion of the IRFR must be affine in the IRFR,
      leading to a Pearson-type diffusion with quadratic volatility.
\item The equilibrium forward curve has an exponential-decay form.
\item The instantaneous short-maturity forward rate equals one-half of 
      the equilibrium regime asymptotic rate at equilibrium.
\item The entire forward curve can be reconstructed from a single 
      forward rate plus structural parameters.
\item The long-time distribution of the IRFR is a shifted 
      inverse-gamma law, heavy-tailed and strictly positive.
\end{itemize}

Taken together, these results yield a unified analytical framework
linking temporal evolution, cross-sectional maturity structure, and
equilibrium regime probabilistic behaviour.

% ------------------------------------------------------------
% 2. Axioms
% ------------------------------------------------------------
\section{Axioms}
\label{sec:axioms}

This section presents the four structural axioms that govern the
dynamics of the instantaneous rolling forward rate. Each axiom has
both an economic interpretation and a mathematical implication, and
together they substantially restrict the class of admissible diffusion
processes governing $x_{t,T}$.

\subsection*{Axiom 1 --- Stochastic Evolution}
\label{ax:stochastic}

At every remaining maturity, the IRFR follows a continuous-state 
diffusion with density $\varphi(x,t,T)$ satisfying the Kolmogorov 
Forward Equation (KFE):
\begin{equation}
\frac{\partial \varphi}{\partial t}
=
-\frac{\partial}{\partial x}(\mu(x,t,T)\varphi)
+ \frac12 \frac{\partial^{2}}{\partial x^{2}}
\!\left(\sigma^{2}(x,t,T)\varphi\right),
\qquad T - t > 0.
\tag{2.0.1}\label{eq:axiom1}
\end{equation}

\subsection*{Axiom 2 --- Rolling--Maturity Invariance}
\label{ax:rolling}

If the drift or diffusion depended directly on calendar time, then two 
IRFRs with the same remaining maturity could exhibit different 
instantaneous dynamics, enabling a trader to roll positions forward and 
exploit systematic differences.

To prevent this form of arbitrage and to ensure dynamic consistency of 
rolling contracts, the IRFR with remaining maturity $\tau = T - t$ must 
satisfy a KFE whose drift and diffusion depend only on $(x,\tau)$:
\begin{equation}
\frac{\partial \varphi}{\partial t}
=
-\frac{\partial}{\partial x}\bigl(\mu(x,\tau)\varphi\bigr)
+ \frac12 \frac{\partial^{2}}{\partial x^{2}}
\bigl(\sigma^{2}(x,\tau)\varphi\bigr),
\qquad \tau > 0.
\tag{2.0.2}\label{eq:axiom2}
\end{equation}

This enforces an invariance structure across calendar time: IRFRs with 
identical remaining maturities must have identical local stochastic 
behaviour.

\subsection*{Axiom 3 --- Equilibrium Moment Transport}
\label{ax:equilibrium}

There exists an equilibrium maturity curve $x_{e,T-t}$ such that, at
equilibrium, the first two moments of the instantaneous rolling forward
rate are transported along constant remaining maturity. In particular,
\begin{equation}
\frac{\partial}{\partial t}\,\mathbb{E}[x_{t,T}]
=
-\frac{\partial}{\partial T}\,\mathbb{E}[x_{t,T}].
\tag{2.0.3}
\end{equation}
and
\begin{equation}
\frac{\partial}{\partial t}\,\mathbb{E}[x_{t,T}^{\,2}]
=
-\frac{\partial}{\partial T}\,\mathbb{E}[x_{t,T}^{\,2}].
\tag{2.0.4}
\end{equation}
Equivalently, the variance satisfies
\begin{equation}
\frac{\partial}{\partial t}\,\mathrm{Var}(x_{t,T})
=
-\frac{\partial}{\partial T}\,\mathrm{Var}(x_{t,T}).
\tag{2.0.5}
\end{equation}
The long-maturity boundary condition is
\[
\lim_{T-t\to\infty} x_{e,T-t} = x_{\infty}.
\]
This axiom imposes equilibrium transport of the mean and variance of the
forward rate across maturity, rather than full stationarity of all
moments.

\subsection*{Axiom 4 --- Positivity at Zero Maturity}
\label{ax:positivity}

At zero time-to-maturity, the IRFR satisfies $x_{t,t}>0$ almost surely. 
To ensure that the boundary at zero is inaccessible, the diffusion 
coefficient must vanish as the short-maturity IRFR approaches zero:
\[
\sigma(x,t,t)\longrightarrow 0
\quad\text{as}\quad x_{t,t}\to 0.
\]
% ------------------------------------------------------------
% 3. Structural Framework
% ------------------------------------------------------------
\section{Structural Framework}
\label{sec:framework}

This section derives the structural implications of the axioms for the
dynamics of the instantaneous rolling forward rate. Beginning from the
Kolmogorov Forward Equation, we derive affine forms for the drift and
diffusion of the IRFR and identify the associated equilibrium relation
between the coefficients. The deterministic
evolution of the expectation, the variance structure, and several key
identities follow naturally.

% ------------------------------------------------------------
\subsection*{3.0 \quad KFE Representation and Notational Setup}
% ------------------------------------------------------------

The Kolmogorov Forward Equation governing the evolution of the IRFR is
\begin{equation}
\frac{\partial \varphi}{\partial t}
=
-\frac{\partial}{\partial x}\!\left[\mu(x,t,T)\varphi\right]
+
\frac{1}{2}\frac{\partial^{2}}{\partial x^{2}}
\!\left[\sigma^{2}(x,t,T)\varphi\right].
\tag{3.0.1}\label{eq:kfe-mu-sigma}
\end{equation}

Define the drift--diffusion pair
\begin{align}
f(x,t,T) &= -\mu(x,t,T), &
g^{2}(x,t,T) &= \tfrac12 \sigma^{2}(x,t,T),
\tag{3.0.2--3.0.3}
\end{align}
so that \eqref{eq:kfe-mu-sigma} becomes
\begin{equation}
\frac{\partial \varphi}{\partial t}
=
\frac{\partial^{2}}{\partial x^{2}}\!\left[g^{2}(x,t,T)\varphi\right]
+
\frac{\partial}{\partial x}\!\left[f(x,t,T)\varphi\right].
\tag{3.0.4}\label{eq:kfe-f-g}
\end{equation}

% ------------------------------------------------------------
\subsection{Affine Forms of Drift and Diffusion}
% ------------------------------------------------------------

\begin{statement}[Affine Forms of Drift and Diffusion]
\label{stmt:affineforms}
\emph{
Assume Axiom~2 (rolling--maturity invariance), so that the drift and
diffusion depend on time only through the remaining maturity
\(\tau:=T-t\):
}
\[
f(x,t,T)=F(x,\tau), \qquad g(x,t,T)=G(x,\tau).
\]
\emph{Assume in addition the following affine-closure hypothesis: for each fixed \(\tau\), the diffusion has state-independent slope,}
\[
\partial_{xx}G(x,\tau)=0,
\]
\emph{
and the net probability flux is affine in the state variable, i.e.
there exist functions \(A(\tau)\) and \(B(\tau)\) such that}
\[
\partial_x\!\bigl(G(x,\tau)^2\bigr)+F(x,\tau)=A(\tau)x+B(\tau).
\]
\emph{
Then both the diffusion and drift are affine in the IRFR state \(x\):
}
\begin{equation}
g(x,t,T)=a(\tau)x+b(\tau),
\qquad
f(x,t,T)=a_1(\tau)x+b_1(\tau),
\qquad \tau=T-t.
\tag{3.1.1}\label{eq:affineforms}
\end{equation}
\end{statement}

\medskip

\noindent\textbf{Derivation.}
Expanding the KFE \eqref{eq:kfe-f-g} gives
\begin{align}
\frac{\partial \varphi}{\partial t}
&=
g^{2}\varphi_{xx}
+\bigl(2(g^{2})_{x}+f\bigr)\varphi_{x}
+\bigl((g^{2})_{xx}+f_{x}\bigr)\varphi
\notag\\
&=
g^{2}\varphi_{xx}
+\bigl(4g g_{x}+f\bigr)\varphi_{x}
+\bigl((g^{2})_{xx}+f_{x}\bigr)\varphi.
\tag{3.1.2}\label{eq:kfe-expanded}
\end{align}
Under Axiom~2 we may write \(f(x,t,T)=F(x,\tau)\) and \(g(x,t,T)=G(x,\tau)\),
where \(\tau=T-t\). The additional affine-closure hypothesis now implies
\[
\partial_{xx}G(x,\tau)=0,
\]
hence for each fixed \(\tau\),
\[
G(x,\tau)=a(\tau)x+b(\tau).
\]
Substituting this into the affine flux identity
\[
\partial_x\!\bigl(G(x,\tau)^2\bigr)+F(x,\tau)=A(\tau)x+B(\tau)
\]
yields
\[
F(x,\tau)
=
A(\tau)x+B(\tau)-\partial_x\!\bigl(G(x,\tau)^2\bigr).
\]
Since \(G(x,\tau)=a(\tau)x+b(\tau)\), we have
\[
\partial_x\!\bigl(G(x,\tau)^2\bigr)
=
2a(\tau)\bigl(a(\tau)x+b(\tau)\bigr),
\]
which is affine in \(x\). Therefore \(F(x,\tau)\) is affine in \(x\) as well:
\[
F(x,\tau)=a_1(\tau)x+b_1(\tau).
\]
Returning to the original notation,
\[
g(x,t,T)=a(T-t)x+b(T-t),
\qquad
f(x,t,T)=a_1(T-t)x+b_1(T-t).
\]
This proves Statement~\ref{stmt:affineforms}.
% ------------------------------------------------------------
\subsection{Equilibrium Relation of Drift and Diffusion}
% ------------------------------------------------------------

\begin{statement}[Equilibrium Relation of Drift and Diffusion]
\label{stmt:eqrelation}
\emph{
Let \(\tau:=T-t\). Suppose the drift and diffusion are affine in the IRFR state:
}
\[
f(x,t,T)=a_1(\tau)x+b_1(\tau),\qquad
g(x,t,T)=a(\tau)x+b(\tau).
\]
\emph{
Assume further that the equilibrium mean depends only on remaining maturity,
}
\[
\mathbb{E}[x_{t,T}]=x_{e,\tau},
\]
\emph{
and that the affine intercepts satisfy the proportionality condition
}
\[
b(\tau)=\frac{a(\tau)}{a_1(\tau)}\,b_1(\tau),
\qquad a_1(\tau)\neq 0.
\]
\emph{
Then the drift and diffusion satisfy
}
\begin{align}
f(x,t,T)
&=
a_1(\tau)\bigl(x-x_{e,\tau}\bigr)+x'_{e,\tau},
\tag{3.2.1a}\label{eq:eqdrift}\\[0.3em]
g(x,t,T)
&=
a(\tau)\bigl(x-x_{e,\tau}\bigr)
+
\frac{a(\tau)}{a_1(\tau)}\,x'_{e,\tau}.
\tag{3.2.1b}\label{eq:eqdiffusion}
\end{align}
\end{statement}
% ------------------------------------------------------------
\subsection{Deterministic Evolution of the IRFR Expectation}
% ------------------------------------------------------------

\begin{statement}[Expected forward-curve dynamics]\label{stmt:expectation}
\normalfont
For fixed maturity \(T\), suppose the expected forward rate evolves according to
\[
\frac{\partial \mathbb{E}[x_{t,T}]}{\partial t}
=
-a_1\bigl(\mathbb{E}[x_{t,T}] - x_{e,T-t}\bigr)
+
\frac{\partial x_{e,T-t}}{\partial t},
\]
where \(a_1>0\) is constant. Then
\begin{equation}
\mathbb{E}[x_{t,T}]
=
x_{e,T-t}
+
\left(x_{t_0,T}-x_{e,T-t_0}\right)e^{-a_1(t-t_0)}.
\tag{3.3.1}\label{eq:expectation-general}
\end{equation}
\end{statement}

Derivation is given in Appendix \ref{appendix:A3}.

% ------------------------------------------------------------
\subsection{Drift Constant Equals Stochastic Constant}
% ------------------------------------------------------------

\begin{statement}[Drift Constant Equals Stochastic Constant]
\label{stmt:a1eqasq}
\emph{
The affine drift constant equals the stochastic constant:
}
\[
a_{1}=a^{2}.
\tag{3.4.1}
\]
\end{statement}

\medskip

\noindent\textbf{Reasoning.}
The identity \(a_{1}=a^{2}\) is a structural equilibrium restriction.
If \(a_{1}\neq a^{2}\), the equilibrium second-moment evolution contains a
residual term that can be eliminated only in the degenerate case of a globally
linear equilibrium curve, which is incompatible with Axioms~3--4. Hence the
only admissible non-degenerate equilibrium is
\[
a_{1}=a^{2}.
\]
A full derivation is given in Appendix~\ref{appendix:A4}.

Interpretively, this is a drift--diffusion relationship: it ties the
drift loading directly to the squared diffusion loading, so that the
transport speed of the mean and the variance-production scale are no
longer independent structural degrees of freedom.

% ------------------------------------------------------------
\subsection{Interim Drift and Diffusion Forms}
% ------------------------------------------------------------

\begin{statement}[Interim Drift and Diffusion]
\label{stmt:interim}
\emph{
If one adopts the baseline exponential equilibrium specification
\[
x_{e,T-t}=x_{\infty}+A e^{-\lambda^{2}(T-t)},
\]
then consistency of \eqref{eq:eqdrift}--\eqref{eq:eqdiffusion} implies
the interim representations
}
\begin{align}
f(x,t,T)
&=
a^{2}x_{t,T}
-
a^{2}x_{\infty}
-
(\lambda^{2}+a^{2})A\, e^{-\lambda^{2}(T-t)},
\tag{3.5.1a}\label{eq:interim-f}\\[0.3em]
g(x,t,T)
&=
a x_{t,T}
-
a x_{\infty}
-
\left(a + \frac{\lambda^{2}}{a}\right)
A\,e^{-\lambda^{2}(T-t)}.
\tag{3.5.1b}\label{eq:interim-g}
\end{align}
\end{statement}

\medskip

This is a conditional consistency result. The exponential equilibrium
curve is imposed as a baseline specification, and \eqref{eq:interim-f}--
\eqref{eq:interim-g} are the corresponding drift and diffusion forms
required by that choice.

% ------------------------------------------------------------
\subsection{Instantaneous Variance}
% ------------------------------------------------------------

\begin{statement}[Instantaneous Variance]
\label{stmt:variance}
\emph{
Under the baseline exponential equilibrium specification used in
Statement~\ref{stmt:interim}, the instantaneous variance evolves according to
}
\begin{equation}
\frac{\partial}{\partial t}\operatorname{Var}(x_{t,T})
=
2\left(
a\,D(t,T)
-
\frac{\lambda^{2}}{a}A e^{-\lambda^{2}(T-t)}
\right)^{2},
\tag{3.6.1}\label{eq:variance}
\end{equation}
where
\begin{equation}
D(t,T)
=
\left(
x_{t_{0},T}
-
x_{\infty}
-
A e^{-\lambda^{2}(T-t_{0})}
\right)
e^{-a^{2}(t-t_{0})}.
\tag{3.6.2}\label{eq:D-def}
\end{equation}
\end{statement}

\medskip

This is the variance evolution implied by the interim forms
\eqref{eq:interim-f}--\eqref{eq:interim-g} under that baseline
specification. A full derivation is given in Appendix~\ref{appendix:A6}.

% ------------------------------------------------------------
\subsection{Market Signal--to--Noise Ratio}
% ------------------------------------------------------------

\begin{statement}[Market Signal--to--Noise Ratio]
\label{stmt:msnr}
\emph{
Define the Market Signal--to--Noise Ratio (MSNR) as the ratio of expected
instantaneous return to the square root of instantaneous variance. Under this
definition, the model-implied MSNR is
}
\[
\mathrm{MSNR}
=
\frac{a}{\sqrt{2}}.
\tag{3.7.1}
\]
\end{statement}

\medskip

In particular, the MSNR is constant: it is independent of calendar time
\(t\), maturity \(T\), and the state \(x_{t,T}\). This does not require
imposing the baseline specification \(\lambda^{2}=a^{2}\).
A derivation is given in Appendix~\ref{appendix:A7}.


% ------------------------------------------------------------
\section{Kolmogorov Forward Equation Solution}
\label{sec:kfe-solution}

This section solves the Kolmogorov Forward Equation for the identified
baseline model. In what follows, we impose
\[
a_{1}=a^{2},
\qquad
\lambda^{2}=a^{2},
\]
as the baseline specification and study the implied equilibrium behaviour of
the IRFR density.

Interpretively, the restriction \(\lambda^{2}=a^{2}\) is a temporal--maturity relationship: it identifies the decay scale in maturity with the
same structural decay scale that governs time evolution in the model.

A convenient shifted coordinate removes the deterministic equilibrium
component and isolates the residual stochastic part of the dynamics. In
those coordinates, the KFE admits a closed-form equilibrium regime given
by a shifted inverse-gamma distribution of order~3.

Throughout, we work with the maturity-indexed equilibrium regime,
parameterised by
\[
\tau:=T-t\ge 0.
\]
Accordingly, distributional results are stated as maturity-indexed
equilibrium densities, while the limit \(\tau\to\infty\) will be
referred to as the \emph{long-maturity limit}.

% ------------------------------------------------------------
\subsection{Maturity-Indexed Equilibrium Density}

Under the baseline specification \(\lambda^{2}=a^{2}\), define the shifted variable
\[
\widetilde{x}
=
x_{t,T}
-
x_{\infty}\!\left(1-e^{-a^{2}(T-t)}\right).
\]

This transformation removes the deterministic equilibrium component and centres the
KFE on the residual stochastic part.

In the shifted coordinate, the forward equation becomes
\begin{equation}
\frac{\partial \varphi}{\partial t}
=
a^{2}\widetilde{x}^{2}\frac{\partial^{2}\varphi}{\partial \widetilde{x}^{2}}
+
5a^{2}\widetilde{x}\frac{\partial \varphi}{\partial \widetilde{x}}
+
3a^{2}\varphi
-
a^{2}x_{\infty}e^{-a^{2}(T-t)}\frac{\partial \varphi}{\partial \widetilde{x}}.
\tag{4.1.1}
\end{equation}

The Frobenius expansion in Appendix~B.1 yields the equilibrium density
\begin{equation}
\varphi(\widetilde{x}, t, T)
=
\frac{x_{\infty}^{3}e^{-3a^{2}(T-t)}}{\Gamma(3)\,\widetilde{x}^{4}}
\exp\!\left(
-\frac{x_{\infty}e^{-a^{2}(T-t)}}{\widetilde{x}}
\right).
\label{eq:shiftedinversegammapdf}
\tag{4.1.2}
\end{equation}

This is a shifted inverse-gamma law with shape parameter \(3\). Its support is strictly
positive, and its heavy right tail reflects the multiplicative structure of the underlying
diffusion. The scale is set by \(x_{\infty}\), while the factor \(e^{-a^{2}(T-t)}\)
controls the maturity-dependent approach to equilibrium.
%================================================================
\subsection{Boundary Behaviour}
\label{sec:boundary}

We now classify the behaviour of the process near the lower boundary
\(x=0\). For the diffusion associated with the IRFR, the scale density is
\begin{equation}
s'(x)
=
\exp\!\left(
-\!\int^{x}
\frac{2f(y,t,T)}{g^{2}(y,t,T)}\,dy
\right)
\sim x^{-2},
\qquad x\to 0^{+},
\tag{4.2.1}\label{eq:scale-density}
\end{equation}
and the corresponding speed density is
\begin{equation}
m(x)
=
\frac{1}{g^{2}(x,t,T)\,s'(x)}
\sim x,
\qquad x\to 0^{+}.
\tag{4.2.2}\label{eq:speed-density}
\end{equation}

Hence
\[
\int_{0^{+}} s'(x)\,dx = \infty,
\qquad
\int_{0^{+}} m(x)\,dx < \infty.
\]
By the standard Feller boundary classification, the origin is therefore
an \emph{entrance boundary}. Equivalently,
\begin{equation}
\boxed{
x=0 \text{ is inaccessible from the interior and cannot be reached in finite time.}
}
\tag{4.2.3}\label{eq:boundary-classification}
\end{equation}

This is fully consistent with Axiom~4 and confirms that the model
preserves positivity of the instantaneous short-maturity IRFR.

% ------------------------------------------------------------
% 5. Structural Integration of Key Results
% ------------------------------------------------------------
\section{Integrated Baseline Specification}
\label{sec:structural-integration}

This section consolidates the structural results established in
Section~3 and Appendix~A. In particular, we now work under the two
baseline identifications
\[
a_{1}=a^{2},
\qquad
\lambda^{2}=a^{2}.
\]
The purpose of the present section is not to re-derive those identities,
but to integrate them into a single coherent specification for the
equilibrium curve, the drift and diffusion coefficients, the Kolmogorov
Forward Equation, and the forward-curve representation.

We proceed in six steps. First, we determine the equilibrium amplitude
\(A\) from the boundary condition. Second, we obtain the fully identified
equilibrium curve. Third, we simplify the drift and diffusion
coefficients. Fourth, we write the associated Kolmogorov Forward
Equation. Fifth, we derive the simplified forward curve representation. Sixth, we state the implied expectation and variance dynamics.

% ------------------------------------------------------------
\subsection{Boundary Condition and Determination of the Amplitude \(A\)}
\label{subsec:determine-A}

From Statement~\ref{stmt:interim}, the interim diffusion is
\[
g(x,t,T)
=
a x_{t,T}
-
a x_{\infty}
-
\left(a+\frac{\lambda^{2}}{a}\right)
A e^{-\lambda^{2}(T-t)}.
\]
Under the baseline specification \(\lambda^{2}=a^{2}\), this becomes
\[
g(x,t,T)
=
a x_{t,T}
-
a x_{\infty}
-
2aA e^{-a^{2}(T-t)}.
\]

Axiom~4 requires that the diffusion vanish at the short-maturity
boundary:
\[
g(x_{\min},t,t)=0.
\]
At the boundary \(x_{t,t}=x_{\min}=0\), and at \(T=t\) one has
\(e^{-a^{2}(T-t)}=1\). Hence
\[
0
=
-a x_{\infty}
-
2aA.
\]
Since \(a\neq 0\), it follows that
\begin{equation}
A=-\frac{x_{\infty}}{2}.
\tag{5.1.1}\label{eq:A-half-xinf}
\end{equation}

Thus the equilibrium amplitude is determined by the boundary
condition.

% ------------------------------------------------------------
\subsection{Fully Identified Equilibrium Curve}
\label{subsec:reevaluate-xe}

Under the baseline specification \(\lambda^{2}=a^{2}\), substituting
\(A=-x_{\infty}/2\) into the equilibrium form
\[
x_{e,T-t}=x_{\infty}+A e^{-\lambda^{2}(T-t)}
\]
gives
\begin{equation}
x_{e,T-t}
=
x_{\infty}
\left(
1-\frac12 e^{-a^{2}(T-t)}
\right).
\tag{5.2.1}\label{eq:xe-final}
\end{equation}

In particular,
\[
x_{e,0}=\frac{x_{\infty}}{2}.
\]
Thus, at equilibrium, the short end of the curve lies exactly halfway
between the lower boundary \(0\) and the long-maturity asymptote
\(x_{\infty}\).

% ------------------------------------------------------------
\subsection{Updated Drift and Diffusion}
\label{subsec:updated-drift-diffusion}

Differentiate \eqref{eq:xe-final} with respect to \(T\):
\[
\frac{\partial x_{e,T-t}}{\partial T}
=
\frac12 x_{\infty} a^{2} e^{-a^{2}(T-t)}.
\]

Now, under the baseline specification \(\lambda^{2}=a^{2}\), substitute
\(A=-x_{\infty}/2\) into the interim forms
\eqref{eq:interim-f}--\eqref{eq:interim-g}. This yields
\begin{align}
g(x_{t,T},t,T)
&=
a x_{t,T}
-
a x_{\infty}\!\left(1-e^{-a^{2}(T-t)}\right),
\tag{5.3.1a}\label{eq:g-updated}
\\[0.4em]
f(x_{t,T},t,T)
&=
a^{2}x_{t,T}
-
a^{2}x_{\infty}\!\left(1-e^{-a^{2}(T-t)}\right).
\tag{5.3.1b}\label{eq:f-updated}
\end{align}

Both coefficients therefore revert toward the same identified
equilibrium level, with diffusion loading \(a\) and drift loading
\(a^{2}\).

% ------------------------------------------------------------
\subsection{Updated Kolmogorov Forward Equation}
\label{subsec:updated-kfe}

Substituting \eqref{eq:g-updated} and \eqref{eq:f-updated} into the
Kolmogorov Forward Equation gives
\begin{align}
\frac{\partial \varphi}{\partial t}
&=
\frac{\partial^{2}}{\partial x^{2}}
\left(
\left[
a x_{t,T}
-
a x_{\infty}\!\left(1-e^{-a^{2}(T-t)}\right)
\right]^{2}
\varphi
\right)
\notag\\[0.4em]
&\qquad
+
\frac{\partial}{\partial x}
\left(
\left[
a^{2}x_{t,T}
-
a^{2}x_{\infty}\!\left(1-e^{-a^{2}(T-t)}\right)
\right]
\varphi
\right).
\tag{5.4.1}\label{eq:kfe-updated}
\end{align}

This is the fully identified forward equation under the baseline
equilibrium specification.

% ------------------------------------------------------------
\subsection{Simplified Forward-Curve Representation}
\label{subsec:revised-forwardcurve}

Using the baseline exponential equilibrium curve together with the baseline forward-curve representation, we obtain
\[
x_{t,T}
=
x_{\infty}
+
A e^{-\lambda^{2}(T-t)}
+
\bigl(x_{t,t}-x_{\infty}-A\bigr)e^{-a^{2}(T-t)}.
\]
Substituting \(\lambda^{2}=a^{2}\) and \(A=-x_{\infty}/2\) gives
\begin{align}
x_{t,T}
&=
x_{\infty}
-
\frac{x_{\infty}}{2}e^{-a^{2}(T-t)}
+
\left(x_{t,t}-x_{\infty}+\frac{x_{\infty}}{2}\right)e^{-a^{2}(T-t)}
\nonumber\\
&=
x_{\infty}
+
(x_{t,t}-x_{\infty})e^{-a^{2}(T-t)}.
\tag{5.5.1}\label{eq:forwardcurve-final}
\end{align}

Thus the entire forward curve is determined by the short-maturity value
\(x_{t,t}\), the equilibrium regime level \(x_{\infty}\), and the common decay rate
\(a^{2}\).

% ------------------------------------------------------------
\subsection{Expectation and Variance under the Identified Baseline}
\label{subsec:instantaneous}

Under the baseline forward-curve form
\eqref{eq:forwardcurve-final}, the expectation is given by
Appendix~\ref{appendix:B3}; see \ref{eq:B3_final_expectation}. Explicitly,
\begin{equation}
\mathbb{E}[x_{t,T}]
=
x_\infty
+
e^{-a^2(T-t_0)}
\left[
\left(x_{t_0,t_0}-\frac{x_\infty}{2}\right)e^{-a^2(t-t_0)}
-
\frac{x_\infty}{2}e^{a^2(t-t_0)}
\right].
\tag{5.6.1}\label{eq:mean-final}
\end{equation}

For the variance, it is safest to retain the direct identified
specialisation of Statement~\ref{stmt:variance}. Using
\(\lambda^{2}=a^{2}\) and \(A=-x_{\infty}/2\) in
\eqref{eq:variance}--\eqref{eq:D-def} gives
\begin{equation}
\frac{\partial}{\partial t}\operatorname{Var}(x_{t,T})
=
2a^{2}
\left(
D(t,T)+\frac{x_{\infty}}{2}e^{-a^{2}(T-t)}
\right)^{2},
\tag{5.6.2}\label{eq:var-final}
\end{equation}
where
\begin{equation}
D(t,T)
=
\left(
x_{t_{0},T}
-
x_{\infty}
+
\frac{x_{\infty}}{2}e^{-a^{2}(T-t_{0})}
\right)e^{-a^{2}(t-t_{0})}.
\tag{5.6.3}\label{eq:D-final}
\end{equation}

Thus the expectation and instantaneous variance both inherit the same
identified decay scale \(a^{2}\), while the cross-sectional shape remains
parametrised by the maturity gap \(T-t\).
% ------------------------------------------------------------
% 6. Economic Interpretation
% ------------------------------------------------------------
\section{Economic Interpretation}
\label{sec:econ}

This section interprets the structural and probabilistic results
obtained in Sections~\ref{sec:framework}--%
\ref{sec:kfe-solution}. Each parameter of the IRFR diffusion plays a
distinct role in shaping both the cross-sectional structure of the
forward curve and its temporal evolution.

% ------------------------------------------------------------
\subsection{Mean-Reversion Dynamics}
\label{subsec:mean-reversion}

The drift and diffusion coefficients in their updated form
\eqref{eq:f-updated}--\eqref{eq:g-updated} may be rewritten as
\begin{align}
f(x,t,T)
&=
a^{2}(x - x_{\infty})
+
a^{2} x_{\infty} e^{-a^{2}(T-t)},
\tag{6.1.1}\label{eq:drift-interpretation}
\\[0.4em]
g(x,t,T)
&=
a(x - x_{\infty})
+
a\, x_{\infty} e^{-a^{2}(T-t)}.
\tag{6.1.2}\label{eq:stochastic-interpretation}
\end{align}

Thus, \(a\) sets the diffusion loading, while \(a^{2}\) governs the
corresponding drift reversion scale.

% ------------------------------------------------------------
\subsection{Equilibrium Regime Level and Equilibrium Structure}
\label{subsec:equilibrium-structure}

The equilibrium curve derived in \eqref{eq:xe-final} is
\begin{equation}
x_{e,T-t}
=
x_{\infty}
-
\frac{x_{\infty}}{2}\, e^{-a^{2}(T-t)},
\tag{6.2.1}\label{eq:xe-equilibrium-interpretation}
\end{equation}
which displays exponential decay from the short end toward the
equilibrium regime asymptote. At zero maturity, equilibrium requires
\[
x_{e,0}=\frac{x_{\infty}}{2},
\]
This expresses the structural link between the short and long ends of
the curve.

% ------------------------------------------------------------
\subsection{Forward-Curve Dynamics}
\label{subsec:forward-curve-dynamics}

Under the baseline specification, the expected forward curve evolves according to
\begin{equation}
\mathbb{E}[x_{t,T}]
=
x_{t_{0},t_{0}}e^{-a^{2}(T+t-2t_{0})}
+
x_{\infty}
-
x_{\infty} e^{-a^{2}(T-t_{0})}.
\tag{6.3.1}\label{eq:forward-curve-dynamics}
\end{equation}
as shown in Appendix~\ref{appendix:B3}; see \eqref{eq:B3_final_expectation}. This expression shows that the mean forward curve is pinned down by the initial short-maturity condition \(x_{t_{0},t_{0}}\), the equilibrium regime level \(x_{\infty}\), and the common decay rate \(a^{2}\). For fixed \(T\), the first term decays with calendar time, while the second and third terms govern convergence toward the equilibrium regime level.

\medskip
\noindent\textbf{Covariance-style structure and equilibrium.}
Introducing shifted variables \(T' = T - t_0\) and \(t' = t - t_0\), so that
\(\tau = T' - t'\), makes the forward-curve dynamics separate into elapsed-time
decay and maturity loading through exponential factors. In this form, the
transition structure is reminiscent of the covariance kernels in classical
affine term-structure models such as Vasicek and CIR: \(t'\) governs temporal
propagation, while \(\tau\) governs maturity dependence.


% ------------------------------------------------------------
\subsection{Distributional Implications}
\label{subsec:distributional}

The equilibrium regime density from Section~4.1, equation \eqref{eq:shiftedinversegammapdf}, is the shifted inverse-gamma law

\begin{equation}
\varphi(\widetilde{x},t,T)
=
\frac{x_{\infty}^{3}e^{-3a^{2}(T-t)}}{\Gamma(3)\,\widetilde{x}^{4}}
\exp\!\left(
-\frac{x_{\infty}e^{-a^{2}(T-t)}}{\widetilde{x}}
\right),
\tag{6.4.1}\label{eq:long-time-limit-density-interpretation}
\end{equation}
which is strictly positive and heavy-tailed. This is consistent with
the empirical observation that interest rates can occasionally exhibit
large moves.

% ------------------------------------------------------------
\subsection{Boundary Behaviour}
\label{subsec:econ-boundary}

As shown in Section~\ref{sec:boundary}, the origin is an entrance
boundary. The IRFR therefore cannot reach zero in finite time from the
interior.
% ------------------------------------------------------------
% 7. Applications
% ------------------------------------------------------------
\section{Applications}
\label{sec:applications}

The structural and probabilistic characterisation of the IRFR provides
a tractable platform for simulation, estimation, curve construction,
risk analysis, and pricing of interest-rate--dependent instruments.
This section outlines several practical uses of the model.

% ------------------------------------------------------------
\subsection{Simulating IRFR Paths}
\label{subsec:simulation}

Using the updated drift and diffusion 
\eqref{eq:f-updated}--\eqref{eq:g-updated}, the IRFR under the physical
measure satisfies the stochastic differential equation
\[
dx_{t,T}
=
-f(x,t,T)\,dt
+
\sqrt{2}g(x,t,T)\,dW_{t}.
\]

Given the explicit forms of $f$ and $g$, one may simulate sample paths
of forward rates using Euler--Maruyama or higher-order schemes. The
affine structure supports stable numerical implementation, while the
boundary properties established in Section~\ref{sec:boundary} preserve
positivity.

% ------------------------------------------------------------
\subsection{Yield Curve Construction}
\label{subsec:yield-curve}

Under the baseline specification, the forward curve can be constructed from
the short rate \(x_{t,t}\) using the representation
\[
x_{t,T}
=
(x_{t,t} - x_{\infty})e^{-a^{2}(T-t)}
+
x_{\infty}.
\]
This exponential form fits naturally into continuous-time curve building
and provides a parsimonious two-parameter representation of the term
structure at any given time.

% ------------------------------------------------------------
\subsection{Risk Management and Stress Testing}
\label{subsec:risk}

Because the maturity-indexed equilibrium density is heavy-tailed and strictly positive,
stress scenarios may incorporate large, sudden rate moves drawn from
the shifted inverse-gamma law. Perturbations in the structural
parameters $(a, x_{\infty})$ then provide a coherent way to stress-test
mean reversion and equilibrium regime rate assumptions.

% ------------------------------------------------------------
\subsection{Pricing Interest-Rate--Dependent Claims}
\label{subsec:pricing}

Under a risk-neutral measure, the drift adjusts to incorporate market 
prices of risk, while the diffusion remains affine. The positivity of 
$x_{t,T}$ renders the model suitable for pricing caps, floors, and 
swaptions, and the exponential structure of the forward curve gives 
closed-form or semi-closed-form solutions for many related derivatives.
% ------------------------------------------------------------
% Appendices
% ------------------------------------------------------------
\begin{appendices}
\section{Technical Derivations for Section 3}
\label{appendix:A}

This appendix provides full derivations of the structural results in
Section~\ref{sec:framework}. All proofs follow a clear step-by-step
format.

% ============================================================
\subsection{Proof of Statement~\ref{stmt:affineforms}}
\label{appendix:A1}
\medskip
\noindent\textbf{Step 1: Expand the KFE.}  
From \eqref{eq:kfe-f-g},
\begin{equation}
\frac{\partial\varphi}{\partial t}
=
g^{2}\varphi_{xx}
+
\bigl(2(g^{2})_{x}+f\bigr)\varphi_{x}
+
\bigl((g^{2})_{xx}+f_{x}\bigr)\varphi.
\tag{A.1.1}
\end{equation}
Equivalently, since \((g^{2})_{x}=2g g_{x}\),
\[
\frac{\partial\varphi}{\partial t}
=
g^{2}\varphi_{xx}
+
(4g g_{x}+f)\varphi_{x}
+
\bigl((g^{2})_{xx}+f_{x}\bigr)\varphi.
\]

\medskip
\noindent\textbf{Step 2: Apply rolling--maturity invariance (Axiom 2).}  
Axiom~2 implies that the coefficients depend on calendar time only through
remaining maturity \(\tau=T-t\). Hence there exist functions \(F\) and
\(G\) such that
\begin{equation}
f(x,t,T)=F(x,\tau),
\qquad
g(x,t,T)=G(x,\tau),
\qquad \tau=T-t.
\tag{A.1.2}
\end{equation}

\medskip
\noindent\textbf{Step 3: Impose the affine-closure hypothesis.}  
Assume that for each fixed \(\tau\),
\begin{equation}
\partial_{xx}G(x,\tau)=0,
\tag{A.1.3}\label{eq:A13}
\end{equation}
and that the net probability flux is affine in \(x\):
\begin{equation}
\partial_x\!\bigl(G(x,\tau)^2\bigr)+F(x,\tau)=A(\tau)x+B(\tau)
\tag{A.1.4}\label{eq:A14}
\end{equation}
for some functions \(A(\tau)\) and \(B(\tau)\).

\medskip
\noindent\textbf{Step 4: Solve for the diffusion.}  
Equation \eqref{eq:A13} implies that \(G\) is affine in \(x\). Therefore
there exist functions \(a(\tau)\) and \(b(\tau)\) such that
\begin{equation}
G(x,\tau)=a(\tau)x+b(\tau).
\tag{A.1.5}\label{eq:A15}
\end{equation}

\medskip
\noindent\textbf{Step 5: Solve for the drift.}  
Substituting \eqref{eq:A15} into \eqref{eq:A14} gives
\[
F(x,\tau)
=
A(\tau)x+B(\tau)-\partial_x\!\bigl(G(x,\tau)^2\bigr).
\]
But
\[
\partial_x\!\bigl(G(x,\tau)^2\bigr)
=
2a(\tau)\bigl(a(\tau)x+b(\tau)\bigr),
\]
which is affine in \(x\). Hence \(F\) is affine in \(x\) as well, so there
exist functions \(a_1(\tau)\) and \(b_1(\tau)\) such that
\begin{equation}
F(x,\tau)=a_1(\tau)x+b_1(\tau).
\tag{A.1.6}\label{eq:A16}
\end{equation}

\medskip
\noindent\textbf{Step 6: Return to the original notation.}  
Using \(\tau=T-t\), equations \eqref{eq:A15} and \eqref{eq:A16} become
\[
g(x,t,T)=a(T-t)x+b(T-t),
\qquad
f(x,t,T)=a_1(T-t)x+b_1(T-t).
\]
Thus both diffusion and drift are affine in the IRFR state \(x\), which
establishes Statement~\ref{stmt:affineforms}.
% ============================================================
% ============================================================
\subsection{Proof of Statement \ref{stmt:eqrelation}}
\label{appendix:A2}

\textbf{Goal.}  
Derive the equilibrium drift--diffusion relations
\eqref{eq:eqdrift}--\eqref{eq:eqdiffusion}.

\medskip
\noindent\textbf{Step 1: Equilibrium identity.}  
Under Axiom 3, the equilibrium IRFR depends only on time-to-maturity, so that
\[
x_{e,T-t}=x_{e,\tau}, \qquad \tau:=T-t.
\]
Hence
\[
\frac{\partial x_{e,T-t}}{\partial t}
=
-\,x'_{e,\tau}.
\tag{A.2.12}
\]

\medskip
\noindent\textbf{Step 2: First-moment evolution.}  
From the Fokker--Planck equation, multiplying by \(x\), integrating over
\(x\in[x_{\min},\infty)\), and assuming the boundary terms vanish, gives
\[
\frac{\partial}{\partial t}\mathbb{E}[x_{t,T}]
=
-\mathbb{E}[f(x,t,T)].
\tag{A.2.13}
\]
Using the affine drift form from Statement~\ref{stmt:affineforms},
\[
f(x,t,T)=a_1(\tau)x+b_1(\tau),
\tag{A.2.14}\label{eq:A214}
\]
it follows that
\[
\frac{\partial}{\partial t}\mathbb{E}[x_{t,T}]
=
-\,a_1(\tau)\mathbb{E}[x_{t,T}] - b_1(\tau).
\tag{A.2.15}\label{eq:A215}
\]

\medskip
\noindent\textbf{Step 3: Impose equilibrium.}  
At equilibrium,
\[
\mathbb{E}[x_{t,T}]=x_{e,\tau},
\]
so \eqref{eq:A215} becomes
\[
\frac{\partial x_{e,\tau}}{\partial t}
=
-\,a_1(\tau)x_{e,\tau}-b_1(\tau).
\tag{A.2.16}
\]
Using \(\partial_t x_{e,\tau}=-x'_{e,\tau}\), we obtain
\[
-\,x'_{e,\tau}
=
-\,a_1(\tau)x_{e,\tau}-b_1(\tau),
\]
hence
\[
b_1(\tau)
=
-\,a_1(\tau)x_{e,\tau}+x'_{e,\tau}.
\tag{A.2.17}\label{eq:A217}
\]

\medskip
\noindent\textbf{Step 4: Drift relation.}  
Substituting \eqref{eq:A217} into \eqref{eq:A214} yields
\[
f(x,t,T)
=
a_1(\tau)x-a_1(\tau)x_{e,\tau}+x'_{e,\tau},
\]
that is,
\[
f(x,t,T)
=
a_1(\tau)\bigl(x-x_{e,\tau}\bigr)+x'_{e,\tau}.
\tag{A.2.18}
\]
Identifying \(x\) with the state variable \(x_{t,T}\) gives
\eqref{eq:eqdrift}.

\medskip
\noindent\textbf{Step 5: Diffusion relation.}  
Using the affine diffusion form from Statement~\ref{stmt:affineforms},
\[
g(x,t,T)=a(\tau)x+b(\tau),
\tag{A.2.19}
\]
together with the proportionality hypothesis
\[
b(\tau)=\frac{a(\tau)}{a_1(\tau)}\,b_1(\tau),
\tag{A.2.20}
\]
and substituting \eqref{eq:A217}, one finds
\[
b(\tau)
=
\frac{a(\tau)}{a_1(\tau)}
\Bigl(-\,a_1(\tau)x_{e,\tau}+x'_{e,\tau}\Bigr)
=
-\,a(\tau)x_{e,\tau}+\frac{a(\tau)}{a_1(\tau)}x'_{e,\tau}.
\tag{A.2.21}
\]
Therefore
\[
g(x,t,T)
=
a(\tau)x-a(\tau)x_{e,\tau}+\frac{a(\tau)}{a_1(\tau)}x'_{e,\tau},
\]
i.e.
\[
g(x,t,T)
=
a(\tau)\bigl(x-x_{e,\tau}\bigr)
+
\frac{a(\tau)}{a_1(\tau)}x'_{e,\tau}.
\tag{A.2.22}
\]
Identifying \(x\) with the state variable \(x_{t,T}\) gives
\eqref{eq:eqdiffusion}.

\medskip
This proves Statement~\ref{stmt:eqrelation}.
% ============================================================
\subsection{Proof of Statement \ref{stmt:expectation}}
\label{appendix:A3}

\textbf{Goal.}
Show that, for fixed maturity \(T\),
\[
\mathbb{E}[x_{t,T}]
=
x_{e,T-t}
+
\left(x_{t_0,T}-x_{e,T-t_0}\right)e^{-a_1(t-t_0)}.
\]

\medskip
\noindent\textbf{Step 1: Start from the expectation dynamics.}
By Statement \ref{stmt:expectation}, the expected forward rate satisfies
\[
\frac{\partial \mathbb{E}[x_{t,T}]}{\partial t}
=
-a_1\bigl(\mathbb{E}[x_{t,T}] - x_{e,T-t}\bigr)
+
\frac{\partial x_{e,T-t}}{\partial t},
\]
where \(a_1>0\) is constant.

\medskip
\noindent\textbf{Step 2: Define the deviation from equilibrium.}
Let
\[
z(t):=\mathbb{E}[x_{t,T}] - x_{e,T-t}.
\]
Then
\[
\frac{dz(t)}{dt}
=
\frac{\partial \mathbb{E}[x_{t,T}]}{\partial t}
-
\frac{\partial x_{e,T-t}}{\partial t}.
\]
Substituting the expectation dynamics gives
\[
\frac{dz(t)}{dt}
=
-a_1\bigl(\mathbb{E}[x_{t,T}] - x_{e,T-t}\bigr)
=
-a_1 z(t).
\]

\medskip
\noindent\textbf{Step 3: Solve the differential equation.}
The solution is
\[
z(t)=z(t_0)e^{-a_1(t-t_0)}.
\]
Hence
\[
\mathbb{E}[x_{t,T}] - x_{e,T-t}
=
\left(\mathbb{E}[x_{t_0,T}] - x_{e,T-t_0}\right)e^{-a_1(t-t_0)}.
\]

\medskip
\noindent\textbf{Step 4: Apply the initial condition.}
At time \(t_0\), the forward rate is known, so
\[
\mathbb{E}[x_{t_0,T}] = x_{t_0,T}.
\]
Therefore
\[
\mathbb{E}[x_{t,T}] - x_{e,T-t}
=
\left(x_{t_0,T} - x_{e,T-t_0}\right)e^{-a_1(t-t_0)}.
\]
Rearranging yields
\begin{equation}
\mathbb{E}[x_{t,T}]
=
x_{e,T-t}
+
\left(x_{t_0,T}-x_{e,T-t_0}\right)e^{-a_1(t-t_0)}.
\tag{A.3.1}\label{eq:A31}
\end{equation}

This proves Statement \ref{stmt:expectation}.
% ============================================================
% ============================================================
\subsection{Proof of Statement~\ref{stmt:a1eqasq}}
\label{appendix:A4}

\textbf{Goal.}
Show that equilibrium transport consistency requires
\[
a_1=a^2.
\]

\medskip
\noindent\textbf{Step 1: Start from the variance dynamics.}
Let
\[
V(t,T):=\operatorname{Var}(x_{t,T}).
\]
From the affine diffusion specification, the forward rate satisfies
\[
dx_{t,T}
=
\bigl(-a_1x_{t,T}+a_1x_{e,T-t}+\partial_t x_{e,T-t}\bigr)\,dt
+
\bigl(a x_{t,T}+b(T-t)\bigr)\,dW_t,
\]
where, by Statement~\ref{stmt:eqrelation},
\[
b(T-t)=-a\,x_{e,T-t}+\frac{1}{a}x'_{e,T-t}.
\]

\medskip
\noindent\textbf{Step 2: Derive the second-moment equation.}
Applying It\^o's formula to \(x_{t,T}^2\) and taking expectations gives
\begin{align}
\frac{\partial}{\partial t}\mathbb{E}[x_{t,T}^2]
&=
2\,\mathbb{E}\!\left[x_{t,T}
\bigl(-a_1x_{t,T}+a_1x_{e,T-t}+\partial_t x_{e,T-t}\bigr)\right]
+
\mathbb{E}\!\left[\bigl(a x_{t,T}+b(T-t)\bigr)^2\right]
\nonumber\\
&=
(-2a_1+a^2)\mathbb{E}[x_{t,T}^2]
+
\bigl(2a_1x_{e,T-t}+2\partial_t x_{e,T-t}+2ab(T-t)\bigr)\mathbb{E}[x_{t,T}]
+
b(T-t)^2.
\tag{A.4.1}\label{eq:A41}
\end{align}

\medskip
\noindent\textbf{Step 3: Pass from the second moment to the variance.}
Since
\[
V(t,T)=\mathbb{E}[x_{t,T}^2]-\mathbb{E}[x_{t,T}]^2,
\]
differentiate with respect to \(t\):
\[
\frac{\partial V}{\partial t}
=
\frac{\partial}{\partial t}\mathbb{E}[x_{t,T}^2]
-
2\,\mathbb{E}[x_{t,T}]
\frac{\partial}{\partial t}\mathbb{E}[x_{t,T}].
\]
Using Statement~\ref{stmt:expectation},
\[
\frac{\partial}{\partial t}\mathbb{E}[x_{t,T}]
=
-a_1\bigl(\mathbb{E}[x_{t,T}]-x_{e,T-t}\bigr)
+
\partial_t x_{e,T-t},
\tag{A.4.2}\label{eq:A42}
\]
and substituting \eqref{eq:A41} and \eqref{eq:A42}, the mixed terms cancel. This yields
\[
\frac{\partial V}{\partial t}
=
2(a^2-a_1)V
+
\bigl(a\,x_{e,T-t}+b(T-t)\bigr)^2.
\tag{A.4.3}\label{eq:A43}
\]

\medskip
\noindent\textbf{Step 4: Use the equilibrium diffusion relation.}
From Statement~\ref{stmt:eqrelation},
\[
b(T-t)=-a\,x_{e,T-t}+\frac{1}{a}x'_{e,T-t}.
\]
Hence
\[
a\,x_{e,T-t}+b(T-t)=\frac{1}{a}x'_{e,T-t},
\]
so \eqref{eq:A43} becomes
\[
\frac{\partial V}{\partial t}
=
2(a^2-a_1)V
+
\frac{1}{a^2}\bigl(x'_{e,T-t}\bigr)^2.
\tag{A.4.4}\label{eq:A44}
\]

\medskip
\noindent\textbf{Step 5: Impose equilibrium transport consistency.}
Under equilibrium transport, admissible term-structure objects depend on calendar time and maturity only through
\[
\tau:=T-t.
\]
In particular, the variance must satisfy
\[
V(t,T)=\widetilde V(\tau)
\]
for some function \(\widetilde V\). Therefore
\[
\frac{\partial V}{\partial t}=-\widetilde V'(\tau),
\]
and \eqref{eq:A44} becomes
\[
-\widetilde V'(\tau)
=
2(a^2-a_1)\widetilde V(\tau)
+
\frac{1}{a^2}\bigl(x_e'(\tau)\bigr)^2.
\tag{A.4.5}\label{eq:A45}
\]

Now suppose \(a_1\neq a^2\). Then the coefficient on \(\widetilde V(\tau)\) in \eqref{eq:A45} is nonzero, so the variance dynamics contain an additional exponential decay or growth term that is independent of the equilibrium transport speed \(a\). This breaks the required one-parameter transport structure: the evolution of the variance would no longer be governed solely by the same decay scale that determines the equilibrium forward-curve propagation.

Therefore the transport-consistent equilibrium representation is possible only if
\[
2(a^2-a_1)=0.
\]
Hence
\[
a_1=a^2.
\tag{A.4.6}\label{eq:A46}
\]

This proves Statement~\ref{stmt:a1eqasq}.

% ============================================================
\subsection{Proof of Statement~\ref{stmt:interim}}
\label{appendix:A5}

\textbf{Goal.}
Show that once the baseline exponential equilibrium specification is
imposed, consistency of the equilibrium relations determines the interim
forms \eqref{eq:interim-f}--\eqref{eq:interim-g}.

\medskip
\noindent\textbf{Step 1: Start from the equilibrium drift and diffusion relations.}

From Appendix~\ref{appendix:A2}, the equilibrium relations are
\begin{align}
f(x,t,T)
&=
a_1 x_{t,T}
-
a_1 x_{e,T-t}
-
\frac{\partial x_{e,T-t}}{\partial t},
\tag{A.5.1a}\label{eq:A5eqdrift}
\\[0.4em]
g(x,t,T)
&=
a x_{t,T}
-
a x_{e,T-t}
-
\frac{a}{a_1}\frac{\partial x_{e,T-t}}{\partial t}.
\tag{A.5.1b}\label{eq:A5eqdiff}
\end{align}
Since \(\tau=T-t\), we have
\[
-\frac{\partial x_{e,T-t}}{\partial t}
=
\frac{\partial x_{e,T-t}}{\partial T}.
\]
Using Statement~\ref{stmt:a1eqasq}, namely \(a_1=a^2\), this becomes
\begin{align}
f(x,t,T)
&=
a^{2}x_{t,T}
-
a^{2}x_{e,T-t}
+
\frac{\partial x_{e,T-t}}{\partial T},
\tag{A.5.2a}\label{eq:A5fpre}
\\[0.4em]
g(x,t,T)
&=
a x_{t,T}
-
a x_{e,T-t}
+
\frac{1}{a}\frac{\partial x_{e,T-t}}{\partial T}.
\tag{A.5.2b}\label{eq:A5gpre}
\end{align}

\medskip
\noindent\textbf{Step 2: Impose the baseline exponential equilibrium specification.}

At this stage, no unique functional form for \(x_{e,T-t}\) has been
derived. We now impose the baseline exponential equilibrium specification
used in the main text:
\begin{equation}
x_{e,T-t}
=
x_{\infty}+A e^{-\lambda^{2}(T-t)}.
\tag{A.5.3}\label{eq:A5xeTT}
\end{equation}
Equivalently, in the maturity variable \(\tau=T-t\),
\begin{equation}
x_{e,\tau}=x_{\infty}+A e^{-\lambda^{2}\tau}.
\tag{A.5.4}\label{eq:A5xe}
\end{equation}
Differentiating \eqref{eq:A5xeTT} with respect to \(T\) yields
\begin{equation}
\frac{\partial x_{e,T-t}}{\partial T}
=
-\lambda^{2}A e^{-\lambda^{2}(T-t)}.
\tag{A.5.5}\label{eq:A5xediff}
\end{equation}

\medskip
\noindent\textbf{Step 3: Substitute into the equilibrium relations.}

Substituting \eqref{eq:A5xeTT} and \eqref{eq:A5xediff} into
\eqref{eq:A5fpre} gives
\begin{align}
f(x,t,T)
&=
a^{2}x_{t,T}
-
a^{2}\!\left(x_{\infty}+A e^{-\lambda^{2}(T-t)}\right)
-
\lambda^{2}A e^{-\lambda^{2}(T-t)}
\nonumber\\
&=
a^{2}x_{t,T}
-
a^{2}x_{\infty}
-
(\lambda^{2}+a^{2})A e^{-\lambda^{2}(T-t)}.
\tag{A.5.6}\label{eq:A51}
\end{align}
Likewise, substituting \eqref{eq:A5xeTT} and \eqref{eq:A5xediff} into
\eqref{eq:A5gpre} gives
\begin{align}
g(x,t,T)
&=
a x_{t,T}
-
a\!\left(x_{\infty}+A e^{-\lambda^{2}(T-t)}\right)
-
\frac{\lambda^{2}}{a}A e^{-\lambda^{2}(T-t)}
\nonumber\\
&=
a x_{t,T}
-
a x_{\infty}
-
\left(a+\frac{\lambda^{2}}{a}\right)
A e^{-\lambda^{2}(T-t)}.
\tag{A.5.7}\label{eq:A52}
\end{align}

\medskip
\noindent\textbf{Conclusion.}

Thus, once the baseline exponential equilibrium curve is imposed,
consistency of the equilibrium drift and diffusion relations yields the
interim forms \eqref{eq:interim-f}--\eqref{eq:interim-g}. This is a
conditional identification result, not a uniqueness statement about all
admissible equilibrium curves.

% ============================================================
\subsection{Proof of Statement~\ref{stmt:variance}}
\label{appendix:A6}

\textbf{Goal.}
Derive the variance evolution implied by the interim drift and diffusion
forms obtained under the baseline exponential equilibrium specification.

\medskip
\noindent\textbf{Step 1: Start from the interim diffusion form.}

From Appendix~\ref{appendix:A5}, once the baseline exponential
equilibrium curve is imposed, the diffusion coefficient is
\begin{equation}
g(x,t,T)
=
a x_{t,T}
-
a x_{\infty}
-
\left(a+\frac{\lambda^{2}}{a}\right)
A e^{-\lambda^{2}(T-t)}.
\tag{A.6.1}\label{eq:A61}
\end{equation}

\medskip
\noindent\textbf{Step 2: Decompose the state into mean plus fluctuation.}

Let
\[
m(t,T):=\mathbb{E}[x_{t,T}].
\]
Under the same baseline specification, the mean forward curve has the
representation
\begin{equation}
m(t,T)
=
x_{\infty}
+
A e^{-\lambda^{2}(T-t)}
+
D(t,T),
\tag{A.6.2}\label{eq:A62}
\end{equation}
where
\begin{equation}
D(t,T)
=
\left(
x_{t_0,T}
-
x_{\infty}
-
A e^{-\lambda^{2}(T-t_0)}
\right)e^{-a^{2}(t-t_0)}.
\tag{A.6.3}\label{eq:A63}
\end{equation}

Now define the centred fluctuation
\[
\widetilde{x}_{t,T}:=x_{t,T}-m(t,T),
\qquad
\mathbb{E}[\widetilde{x}_{t,T}]=0.
\]
Substituting \(x_{t,T}=m(t,T)+\widetilde{x}_{t,T}\) together with
\eqref{eq:A62} into \eqref{eq:A61} gives
\begin{align}
g(x,t,T)
&=
a\bigl(m(t,T)+\widetilde{x}_{t,T}\bigr)
-
a x_{\infty}
-
\left(a+\frac{\lambda^{2}}{a}\right)A e^{-\lambda^{2}(T-t)}
\nonumber\\
&=
a\widetilde{x}_{t,T}
+
aD(t,T)
-
\frac{\lambda^{2}}{a}A e^{-\lambda^{2}(T-t)}.
\tag{A.6.4}\label{eq:A64}
\end{align}

\medskip
\noindent\textbf{Step 3: Apply the second-moment identity.}

The second moment satisfies
\begin{equation}
\frac{\partial}{\partial t}\mathbb{E}[x_{t,T}^{2}]
=
-2\,\mathbb{E}[x_{t,T}f(x,t,T)]
+
2\,\mathbb{E}[g(x,t,T)^{2}].
\tag{A.6.5}\label{eq:A65}
\end{equation}

Using
\[
\operatorname{Var}(x_{t,T})
=
\mathbb{E}[x_{t,T}^{2}] - m(t,T)^{2},
\]
we obtain
\begin{align}
\frac{\partial}{\partial t}\operatorname{Var}(x_{t,T})
&=
\frac{\partial}{\partial t}\mathbb{E}[x_{t,T}^{2}]
-
\frac{\partial}{\partial t}m(t,T)^{2}
\nonumber\\
&=
-2\Bigl(\mathbb{E}[x_{t,T}f(x,t,T)]-m(t,T)\mathbb{E}[f(x,t,T)]\Bigr)
+
2\,\mathbb{E}[g(x,t,T)^{2}].
\tag{A.6.6}\label{eq:A66}
\end{align}

From Appendix~\ref{appendix:A5}, the drift is affine:
\[
f(x,t,T)
=
a^{2}x_{t,T}
-
a^{2}x_{\infty}
-
(\lambda^{2}+a^{2})A e^{-\lambda^{2}(T-t)}.
\]
Therefore its covariance with \(x_{t,T}\) is simply
\[
\mathbb{E}[x_{t,T}f(x,t,T)]-m(t,T)\mathbb{E}[f(x,t,T)]
=
a^{2}\operatorname{Var}(x_{t,T}),
\]
and \eqref{eq:A66} becomes
\begin{equation}
\frac{\partial}{\partial t}\operatorname{Var}(x_{t,T})
=
-2a^{2}\operatorname{Var}(x_{t,T})
+
2\,\mathbb{E}[g(x,t,T)^{2}].
\tag{A.6.7}\label{eq:A67}
\end{equation}

\medskip
\noindent\textbf{Step 4: Evaluate the diffusion-square expectation.}

From \eqref{eq:A64}, write
\[
g(x,t,T)
=
a\widetilde{x}_{t,T}
+
b(t,T),
\qquad
b(t,T):=
aD(t,T)-\frac{\lambda^{2}}{a}A e^{-\lambda^{2}(T-t)}.
\]
Hence
\begin{align}
\mathbb{E}[g(x,t,T)^{2}]
&=
\mathbb{E}\!\left[\left(a\widetilde{x}_{t,T}+b(t,T)\right)^{2}\right]
\nonumber\\
&=
a^{2}\mathbb{E}[\widetilde{x}_{t,T}^{2}]
+
2ab(t,T)\mathbb{E}[\widetilde{x}_{t,T}]
+
b(t,T)^{2}
\nonumber\\
&=
a^{2}\operatorname{Var}(x_{t,T})
+
\left(
aD(t,T)-\frac{\lambda^{2}}{a}A e^{-\lambda^{2}(T-t)}
\right)^{2},
\tag{A.6.8}\label{eq:A68}
\end{align}
because \(\mathbb{E}[\widetilde{x}_{t,T}]=0\).

Substituting \eqref{eq:A68} into \eqref{eq:A67} yields
\begin{align}
\frac{\partial}{\partial t}\operatorname{Var}(x_{t,T})
&=
-2a^{2}\operatorname{Var}(x_{t,T})
+
2a^{2}\operatorname{Var}(x_{t,T})
\nonumber\\
&\quad
+
2\left(
aD(t,T)-\frac{\lambda^{2}}{a}A e^{-\lambda^{2}(T-t)}
\right)^{2}
\nonumber\\
&=
2\left(
aD(t,T)-\frac{\lambda^{2}}{a}A e^{-\lambda^{2}(T-t)}
\right)^{2}.
\tag{A.6.9}\label{eq:A69}
\end{align}

\medskip
\noindent\textbf{Conclusion.}

Thus, under the baseline exponential equilibrium specification, the
variance evolution implied by the interim affine forms is exactly
\[
\frac{\partial}{\partial t}\operatorname{Var}(x_{t,T})
=
2\left(
aD(t,T)-\frac{\lambda^{2}}{a}A e^{-\lambda^{2}(T-t)}
\right)^{2},
\]
which proves Statement~\ref{stmt:variance}.

% ============================================================
\subsection{Proof of Statement~\ref{stmt:msnr}}
\label{appendix:A7}

\textbf{Goal.}
Derive the market signal-to-noise ratio.

\medskip
\noindent\textbf{Step 1: Start from the variance rate.}

From Appendix~\ref{appendix:A6}, the instantaneous variance rate is
\begin{equation}
\frac{\partial}{\partial t}\operatorname{Var}(x_{t,T})
=
2\left(
aD(t,T)-\frac{\lambda^{2}}{a}A e^{-\lambda^{2}(T-t)}
\right)^{2}.
\tag{A.7.1}\label{eq:A71}
\end{equation}
Hence the instantaneous noise scale is
\begin{equation}
\sqrt{\frac{\partial}{\partial t}\operatorname{Var}(x_{t,T})}
=
\sqrt{2}\,
\left|
aD(t,T)-\frac{\lambda^{2}}{a}A e^{-\lambda^{2}(T-t)}
\right|.
\tag{A.7.2}\label{eq:A72}
\end{equation}

\medskip
\noindent\textbf{Step 2: Evaluate the drift on the mean branch.}

From Appendix~\ref{appendix:A5}, the interim drift is
\begin{equation}
f(x,t,T)
=
a^{2}x_{t,T}
-
a^{2}x_{\infty}
-
(\lambda^{2}+a^{2})A e^{-\lambda^{2}(T-t)}.
\tag{A.7.3}\label{eq:A73}
\end{equation}

Let
\[
m(t,T):=\mathbb{E}[x_{t,T}]
=
x_{\infty}+A e^{-\lambda^{2}(T-t)}+D(t,T),
\]
as obtained in Statement~\ref{stmt:variance}. Evaluating the drift along
this mean branch gives
\begin{align}
f\bigl(m(t,T),t,T\bigr)
&=
a^{2}\Bigl(x_{\infty}+A e^{-\lambda^{2}(T-t)}+D(t,T)\Bigr)
-
a^{2}x_{\infty}
-
(\lambda^{2}+a^{2})A e^{-\lambda^{2}(T-t)}
\nonumber\\
&=
a^{2}D(t,T)-\lambda^{2}A e^{-\lambda^{2}(T-t)}
\nonumber\\
&=
a\left(
aD(t,T)-\frac{\lambda^{2}}{a}A e^{-\lambda^{2}(T-t)}
\right).
\tag{A.7.4}\label{eq:A74}
\end{align}

Therefore the instantaneous signal scale, defined by the magnitude of the
drift along the mean branch, is
\begin{equation}
\left|f\bigl(m(t,T),t,T\bigr)\right|
=
a\left|
aD(t,T)-\frac{\lambda^{2}}{a}A e^{-\lambda^{2}(T-t)}
\right|.
\tag{A.7.5}\label{eq:A75}
\end{equation}

\medskip
\noindent\textbf{Step 3: Form the ratio.}

Hence
\begin{align}
\mathrm{MSNR}
&=
\frac{\left|f\bigl(m(t,T),t,T\bigr)\right|}
{\sqrt{\partial_t\operatorname{Var}(x_{t,T})}}
\nonumber\\
&=
\frac{
a\left|
aD(t,T)-\frac{\lambda^{2}}{a}A e^{-\lambda^{2}(T-t)}
\right|
}{
\sqrt{2}\left|
aD(t,T)-\frac{\lambda^{2}}{a}A e^{-\lambda^{2}(T-t)}
\right|
}
\nonumber\\
&=
\frac{a}{\sqrt{2}}.
\tag{A.7.6}\label{eq:A76}
\end{align}

\medskip
\noindent\textbf{Conclusion.}

Therefore,
\[
\mathrm{MSNR}=\frac{a}{\sqrt{2}},
\]
which proves Statement~\ref{stmt:msnr}.

% ============================================================



% ============================================================
% Appendix B
%==============================================================

\section{Solution to the KFE in the equilibrium regime}
\label{appendix:B}

% ------------------------------------------------------------
\subsection{Full Frobenius Proof of the equilibrium regime PDF}
\label{appendix:B1}

We begin from the transformed KFE in the shifted variable
\[
\widetilde{x}
=
x_{t,T}
-
x_{\infty}\!\left(1-e^{-\lambda^{2}(T-t)}\right),
\]
so that the equilibrium mean has been removed, but without yet imposing any
relation between \(\lambda\) and \(a\). In these coordinates the KFE takes the form
\begin{equation}
\frac{\partial \varphi}{\partial t}
=
a^{2}\widetilde{x}^{2}\varphi_{\widetilde{x}\widetilde{x}}
+
5a^{2}\widetilde{x}\,\varphi_{\widetilde{x}}
+
3a^{2}\varphi
-
\lambda^{2}x_{\infty}e^{-\lambda^{2}(T-t)}\varphi_{\widetilde{x}}.
\tag{B.1.1}\label{eq:B11}
\end{equation}

To analyse the singular structure at \(\widetilde{x}=0\), introduce the reciprocal
variable
\[
z=\frac{1}{\widetilde{x}},
\qquad
z_{\widetilde{x}}=-z^{2},
\qquad
z_{\widetilde{x}\widetilde{x}}=2z^{3}.
\]
Using
\[
\varphi_{\widetilde{x}}=-z^{2}\varphi_{z},
\qquad
\varphi_{\widetilde{x}\widetilde{x}}
=
z^{4}\varphi_{zz}+2z^{3}\varphi_{z},
\]
equation \eqref{eq:B11} becomes
\begin{equation}
\frac{\partial \varphi}{\partial t}
=
a^{2}z^{2}\varphi_{zz}
-
3a^{2}z\varphi_{z}
+
3a^{2}\varphi
+
\lambda^{2}x_{\infty}e^{-\lambda^{2}(T-t)}z^{2}\varphi_{z}.
\tag{B.1.2}\label{eq:B12}
\end{equation}

\medskip
\noindent\textbf{Step 1: Frobenius expansion.}

Seek a Frobenius series of the form
\[
\varphi(z,t,T)=\sum_{i=0}^{\infty} c_i(t)\,z^{i+m}.
\]
Then
\[
\varphi_z
=
\sum_{i=0}^{\infty}(i+m)c_i(t)\,z^{i+m-1},
\qquad
\varphi_{zz}
=
\sum_{i=0}^{\infty}(i+m)(i+m-1)c_i(t)\,z^{i+m-2}.
\]
Substituting into \eqref{eq:B12} and matching coefficients of \(z^{i+m}\) yields
\begin{equation}
\frac{\partial c_i}{\partial t}
=
a^{2}(i+m-1)(i+m-3)c_i
+
\lambda^{2}x_{\infty}e^{-\lambda^{2}(T-t)}(i+m-1)c_{i-1},
\tag{B.1.3}\label{eq:B13}
\end{equation}
with \(c_{-1}\equiv 0\).

\medskip
\noindent\textbf{Step 2: Separate the time dependence.}

Write
\[
c_i(t)=c_i^{*}\,x_{\infty}^{\,i+n}e^{-(i+n)\lambda^{2}(T-t)}.
\]
Then
\[
\frac{\partial c_i}{\partial t}=(i+n)\lambda^{2}c_i(t).
\]
Substituting into \eqref{eq:B13} and cancelling the common factor
\(x_{\infty}^{\,i+n}e^{-(i+n)\lambda^{2}(T-t)}\) gives
\[
(i+n)\lambda^{2}c_i^{*}
=
a^{2}(i+m-1)(i+m-3)c_i^{*}
+
\lambda^{2}(i+m-1)c_{i-1}^{*}.
\]
At \(i=0\), since \(c_{-1}^{*}=0\), this reduces to the indicial relation
\begin{equation}
n\,\frac{\lambda^{2}}{a^{2}}=(m-1)(m-3).
\tag{B.1.4}\label{eq:B14}
\end{equation}
For \(i\ge 1\), subtracting \eqref{eq:B14} gives
\[
\lambda^{2} i\,c_i^{*}
=
a^{2}i(i+2m-4)c_i^{*}
+
\lambda^{2}(i+m-1)c_{i-1}^{*},
\]
hence
\begin{equation}
\frac{c_i^{*}}{c_{i-1}^{*}}
=
-\frac{\frac{\lambda^{2}}{a^{2}}(i+m-1)}{i\left(i+2m-4-\frac{\lambda^{2}}{a^{2}}\right)}.
\tag{B.1.5}\label{eq:B15}
\end{equation}

\medskip
\noindent\textbf{Step 3: Characterise the non-terminating admissible branch.}

Let
\[
r:=\frac{\lambda^{2}}{a^{2}}.
\]
For the Frobenius series to sum to a single inverse-gamma-type branch,
the coefficient ratio needs to have the form
\[
\frac{c_i^{*}}{c_{i-1}^{*}}=-\frac{\beta}{i}
\]
for some constant \(\beta>0\); equivalently, the additional affine factor
in \eqref{eq:B15} needs to reduce to an \(i\)-independent constant. Since the
numerator and denominator in \eqref{eq:B15} are both affine in \(i\) with
the same slope, this occurs if and only if they are identical, namely
\[
i+m-1=i+2m-4-r.
\]
Therefore
\begin{equation}
m=3+r.
\tag{B.1.6}\label{eq:B16}
\end{equation}
Substituting \eqref{eq:B16} into \eqref{eq:B15} gives
\begin{equation}
\frac{c_i^{*}}{c_{i-1}^{*}}=-\frac{r}{i},
\qquad
c_i^{*}=\frac{(-r)^i}{i!}\,c_0^{*}.
\tag{B.1.7}\label{eq:B17}
\end{equation}
The indicial relation \eqref{eq:B14} then yields
\begin{equation}
n=r+2.
\tag{B.1.8}\label{eq:B18}
\end{equation}
Hence the admissible branch is the one-parameter family
\[
\varphi(z,t,T)
=
c_0^{*}x_{\infty}^{r+2}z^{r+3}e^{-(r+2)\lambda^{2}(T-t)}
\exp\!\bigl[-r x_{\infty}e^{-\lambda^{2}(T-t)}z\bigr].
\]
Returning to \(\widetilde{x}=1/z\), this becomes
\begin{equation}
\varphi(\widetilde{x},t,T)
=
c_0^{*}x_{\infty}^{r+2}e^{-(r+2)\lambda^{2}(T-t)}
\widetilde{x}^{-(r+3)}
\exp\!\left(-\frac{r x_{\infty}e^{-\lambda^{2}(T-t)}}{\widetilde{x}}\right).
\tag{B.1.9}\label{eq:B19}
\end{equation}
Thus the Frobenius construction yields an inverse-gamma family with shape
parameter
\[
\alpha=r+2.
\]

\medskip
\noindent\textbf{Step 4: Select the structural branch.}

The branch retained in the manuscript is the order-3 shifted
inverse-gamma law, equivalently the branch with tail
\(\widetilde{x}^{-4}\) and finite second moment. Since \eqref{eq:B19} has tail exponent \(-(r+3)\), this requires
\[
r+3=4,
\qquad\text{so that}\qquad
r=1.
\]
Therefore
\begin{equation}
\boxed{\lambda^{2}=a^{2}}.
\tag{B.1.10}\label{eq:B110}
\end{equation}
From \eqref{eq:B16} and \eqref{eq:B18} we then obtain
\[
m=4,
\qquad
n=3.
\]
The recurrence \eqref{eq:B17} reduces to
\[
\frac{c_i^{*}}{c_{i-1}^{*}}=-\frac{1}{i},
\qquad
c_i^{*}=\frac{(-1)^i}{i!}\,c_0^{*},
\]
so the Frobenius series sums to
\[
\varphi(z,t,T)
=
c_0^{*}x_{\infty}^{3}z^{4}e^{-3a^{2}(T-t)}
\exp\!\bigl[-x_{\infty}e^{-a^{2}(T-t)}z\bigr].
\]
Returning to \(\widetilde{x}=1/z\) gives
\begin{equation}
\varphi(\widetilde{x},t,T)
=
c_0^{*}
\frac{x_{\infty}^{3}e^{-3a^{2}(T-t)}}{\widetilde{x}^{4}}
\exp\!\left(
-\frac{x_{\infty}e^{-a^{2}(T-t)}}{\widetilde{x}}
\right).
\tag{B.1.11}\label{eq:B111}
\end{equation}
Normalisation fixes \(c_0^{*}=1/\Gamma(3)\), and therefore
\begin{equation}
\boxed{
\varphi(\widetilde{x},t,T)
=
\frac{x_{\infty}^{3}e^{-3a^{2}(T-t)}}{\Gamma(3)\,\widetilde{x}^{4}}
\exp\!\left(
-\frac{x_{\infty}e^{-a^{2}(T-t)}}{\widetilde{x}}
\right)
}
\tag{B.1.12}\label{eq:B112}
\end{equation}
which is the shifted inverse-gamma distribution of order \(3\).

% ------------------------------------------------------------
\subsection{Direct Laplace-Transform Derivation of the MGF}
\label{appendix:B2}

Let
\[
\beta_{t,T}
=
x_{\infty}e^{-a^{2}(T-t)}.
\]

From \eqref{eq:B112}, the shifted variable \(\widetilde{x}\) has density
\[
\varphi(\widetilde{x},t,T)
=
\frac{\beta_{t,T}^{3}}{\Gamma(3)\,\widetilde{x}^{4}}
\exp\!\left(
-\frac{\beta_{t,T}}{\widetilde{x}}
\right),
\qquad
\widetilde{x}>0.
\]

We therefore compute its Laplace transform directly. For \(q>0\), define
\[
\widehat{\varphi}(q,t,T)
=
\int_{0}^{\infty} e^{-q\widetilde{x}}\varphi(\widetilde{x},t,T)\,d\widetilde{x}.
\]

Substituting the density gives
\begin{equation}
\widehat{\varphi}(q,t,T)
=
\frac{\beta_{t,T}^{3}}{\Gamma(3)}
\int_{0}^{\infty}
\widetilde{x}^{-4}
\exp\!\left(
-q\widetilde{x}
-\frac{\beta_{t,T}}{\widetilde{x}}
\right)
d\widetilde{x}.
\tag{B.2.1}
\end{equation}

Using the standard integral identity
\[
\int_{0}^{\infty}
x^{\nu-1}
\exp\!\left(-\mu x-\eta/x\right)\,dx
=
2\left(\frac{\eta}{\mu}\right)^{\nu/2}
K_{\nu}\!\left(2\sqrt{\mu\eta}\right),
\]
with \(\nu=-3\), \(\mu=q\), and \(\eta=\beta_{t,T}\), and using
\(K_{-\nu}=K_{\nu}\), we obtain
\begin{equation}
\widehat{\varphi}(q,t,T)
=
\frac{2\beta_{t,T}^{3}}{\Gamma(3)}
\left(\frac{q}{\beta_{t,T}}\right)^{3/2}
K_{3}\!\left(2\sqrt{\beta_{t,T}q}\right).
\tag{B.2.2}
\end{equation}

Since \(\Gamma(3)=2\), this simplifies to
\begin{equation}
\widehat{\varphi}(q,t,T)
=
(\beta_{t,T}q)^{3/2}
K_{3}\!\left(2\sqrt{\beta_{t,T}q}\right)
=
\left(
x_{\infty}e^{-a^{2}(T-t)}q
\right)^{3/2}
K_{3}\!\left(
2\sqrt{x_{\infty}e^{-a^{2}(T-t)}q}
\right).
\tag{B.2.3}
\end{equation}

If one prefers the moment-generating-function variable \(s\), write \(q=-s\).
Then the transform is well defined for \(s<0\) and becomes
\begin{equation}
\widetilde{\varphi}(s,t,T)
=
\left(
-x_{\infty}e^{-a^{2}(T-t)}s
\right)^{3/2}
K_{3}\!\left(
2\sqrt{-x_{\infty}e^{-a^{2}(T-t)}s}
\right),
\qquad s<0.
\tag{B.2.4}
\end{equation}

This is exactly the Laplace transform, equivalently the MGF evaluated on its
negative argument, of the inverse-gamma law of order~3.
% ------------------------------------------------------------

\subsection{Derivation of the Identified Expectation Formula}
\label{appendix:B3}

\textbf{Goal.}
Derive the closed-form expression for the conditional expectation
$\mathbb{E}[x_{t,T}]$ under the identified baseline specification, using
only the first-moment identity implied by the Kolmogorov Forward Equation.

\medskip
\noindent\textbf{Context.}
Appendix~\ref{appendix:A3} establishes the general expectation dynamics under affine
drift.  In this appendix we derive the explicit closed-form expectation
after imposing the identified baseline relations
$a_1 = \lambda^2 = a^2$ and the corresponding drift specification.

\medskip
\noindent\textbf{Step 1: First-moment identity.}
For fixed maturity $T$, the Kolmogorov Forward Equation implies
\begin{equation}
\frac{\partial}{\partial t}\mathbb{E}[x_{t,T}]
=
-\mathbb{E}\!\left[f(x_{t,T},t,T)\right].
\label{eq:B3_first_moment_identity}
\tag{B.3.1}
\end{equation}

\medskip
\noindent\textbf{Step 2: Identified drift.}
Under the identified baseline structure section (\ref{subsec:updated-drift-diffusion}), the drift is
\begin{equation}
f(x_{t,T},t,T)
=
a^{2}x_{t,T}
-
a^{2}x_{\infty}\!\left(1-e^{-a^{2}(T-t)}\right).
\label{eq:B3_drift}
\tag{B.3.2}
\end{equation}

\medskip
\noindent\textbf{Step 3: Taking expectations.}
Taking expectations of \ref{eq:B3_drift} and using linearity yields
\[
\mathbb{E}[f(x_{t,T},t,T)]
=
a^{2}\mathbb{E}[x_{t,T}]
-
a^{2}x_{\infty}\!\left(1-e^{-a^{2}(T-t)}\right).
\]
Substituting into \eqref{eq:B3_first_moment_identity} gives the evolution
equation
\begin{equation}
\frac{\partial}{\partial t}\mathbb{E}[x_{t,T}]
+
a^{2}\mathbb{E}[x_{t,T}]
=
a^{2}x_{\infty}
-
a^{2}x_{\infty}e^{-a^{2}(T-t)}.
\label{eq:B3_expectation_ode}
\tag{B.3.3}
\end{equation}

\medskip
\noindent\textbf{Step 4: Integrating factor.}
Equation \eqref{eq:B3_expectation_ode} is a first-order linear ODE in
$\mathbb{E}[x_{t,T}]$.  The integrating factor is $e^{a^{2}t}$, and
multiplying \eqref{eq:B3_expectation_ode} by this factor yields
\begin{equation}
\frac{\partial}{\partial t}
\left(
e^{a^{2}t}\mathbb{E}[x_{t,T}]
\right)
=
a^{2}x_{\infty}e^{a^{2}t}
-
a^{2}x_{\infty}e^{a^{2}t}e^{-a^{2}(T-t)}.
\notag
\end{equation}

\medskip
\noindent\textbf{Step 5: Simplifying the forcing term.}
Using
\[
e^{a^{2}t}e^{-a^{2}(T-t)} = e^{-a^{2}T}e^{2a^{2}t},
\]
the equation becomes
\begin{equation}
\frac{\partial}{\partial t}
\left(
e^{a^{2}t}\mathbb{E}[x_{t,T}]
\right)
=
a^{2}x_{\infty}e^{a^{2}t}
-
a^{2}x_{\infty}e^{-a^{2}T}e^{2a^{2}t}.
\notag
\end{equation}

\medskip
\noindent\textbf{Step 6: Integration in time.}
Integrating from $t_0$ to $t$ gives
\begin{align}
e^{a^{2}t}\mathbb{E}[x_{t,T}]
-
e^{a^{2}t_0}\mathbb{E}[x_{t_0,T}]
&=
x_{\infty}\bigl(e^{a^{2}t}-e^{a^{2}t_0}\bigr)
-
\frac{x_{\infty}}{2}e^{-a^{2}T}
\bigl(e^{2a^{2}t}-e^{2a^{2}t_0}\bigr).
\notag
\end{align}

\medskip
\noindent\textbf{Step 7: Solving for the expectation.}
Dividing through by $e^{a^{2}t}$ yields
\begin{align}
\mathbb{E}[x_{t,T}]
&=
\mathbb{E}[x_{t_0,T}]e^{-a^{2}(t-t_0)}
+
x_{\infty}\bigl(1-e^{-a^{2}(t-t_0)}\bigr)
\nonumber\\
&\quad
-
\frac{x_{\infty}}{2}e^{-a^{2}(T-t)}
+
\frac{x_{\infty}}{2}e^{-a^{2}(T+t-2t_0)}.
\label{eq:B3_pre_final}
\tag{B.3.4}
\end{align}

\medskip
\noindent\textbf{Step 8: Imposing the initial condition.}
Using
\[
\mathbb{E}[x_{t_0,T}] = x_{t_0,T}
=
(x_{t_0,t_0}-x_{\infty})e^{-a^{2}(T-t_0)} + x_{\infty},
\]
and simplifying \eqref{eq:B3_pre_final}, we obtain the final exponential
form
\begin{equation}
\boxed{
\mathbb{E}[x_{t,T}]
=
x_{\infty}
+
e^{-a^{2}(T-t_0)}
\left[
\left(x_{t_0,t_0}-\frac{x_{\infty}}{2}\right)e^{-a^{2}(t-t_0)}
-
\frac{x_{\infty}}{2}e^{a^{2}(t-t_0)}
\right].
}
\label{eq:B3_final_expectation}
\tag{B.3.5}
\end{equation}
Hence, the expectation admits the hyperbolic form
\begin{equation}
\boxed{
\mathbb{E}[x_{t,T}]
=
x_\infty
+
e^{-a^2(T-t_0)}
\Big[
(x_{t_0,t_0}-x_\infty)\cosh\!\big(a^2(t-t_0)\big)
-
x_{t_0,t_0}\sinh\!\big(a^2(t-t_0)\big)
\Big].
}
\tag{B.3.6}
\end{equation}
% 

%---------------------------------------
%\subsection{Derivation of the Identified Expectation Formula}
%\label{appendix:B3}

%\textbf{Goal.}
%Derive the correct closed-form expression for the conditional expectation
%\(\mathbb{E}[x_{t,T}]\) and present it in both exponential and hyperbolic form.

%\medskip
%\noindent\textbf{Step 1: Starting point.}
%From equation \((2.3.7)\), the time derivative of the expectation %is
%\begin{equation}
%\frac{\partial \mathbb{E}[x_{t,T}]}{\partial t}
%=
%-a^2
%\left[
%\left(x_{t_0,t_0}-\frac{x_\infty}{2}\right)e^{-a^2(t-t_0)}
%-
%\frac{x_\infty}{2}e^{a^2(t-t_0)}
%\right]
%e^{-a^2(T-t_0)}.
%\tag{B.3.1}
%\end{equation}
%We impose the boundary condition
%\[
%\mathbb{E}[x_{t_0,T}] = x_{t_0,T}.
%\]

%\medskip
%\noindent\textbf{Step 2: Integration in $t$.}
%Integrating from \(t_0\) to \(t\) yields
%\begin{align}
%\mathbb{E}[x_{t,T}] - x_{t_0,T}
%&=
%-a^2 e^{-a^2(T-t_0)}
%\int_{t_0}^{t}
%\left[
%\left(x_{t_0,t_0}-\frac{x_\infty}{2}\right)e^{-a^2(s-t_0)}
%-
%\frac{x_\infty}{2}e^{a^2(s-t_0)}
%\right] ds
%\nonumber\\[4pt]
%\mathbb{E}[x_{t,T}]
%&=
%e^{-a^2(T-t_0)}
%\Bigg[
%\left(x_{t_0,t_0}-\frac{x_\infty}{2}\right)
%\left(e^{-a^2(t-t_0)}-1\right)
%-
%\frac{x_\infty}{2}
%\left(e^{a^2(t-t_0)}-1\right)
%\Bigg].
%\label{eq:expectedvalue}
%\tag{B.3.2}
%\end{align}
%
%\medskip
%\noindent\textbf{Step 3: Substituting the initial curve value.}
%Using
%\[
%x_{t_0,T}
%=
%(x_{t_0,t_0}-x_\infty)e^{-a^2(T-t_0)} + x_\infty,
%\]
%and simplifying, we obtain
%\begin{equation}
%\mathbb{E}[x_{t,T}]
%=
%x_\infty
%+
%e^{-a^2(T-t_0)}
%\left[
%\left(x_{t_0,t_0}-\frac{x_\infty}{2}\right)e^{-a^2(t-t_0)}
%-
%\frac{x_\infty}{2}e^{a^2(t-t_0)}
%\right].
%\label{eq:expectationformuala}
%\tag{B.3.3}
%\end{equation}

%\medskip
%\noindent\textbf{Step 4: Hyperbolic representation.}
%Let \(\tau = a^2(t-t_0)\). Using
%\[
%e^{\pm \tau} = \cosh\tau \pm \sinh\tau,
%\]
%the bracketed term in \ref{eq:expectationformuala} can be %rewritten as
%\[
%(x_{t_0,t_0}-x_\infty)\cosh\tau
%-
%x_{t_0,t_0}\sinh\tau.
%\]
%Hence, the expectation admits the hyperbolic form
%\begin{equation}
%\boxed{
%\mathbb{E}[x_{t,T}]
%=
%x_\infty
%+
%e^{-a^2(T-t_0)}
%\Big[
%(x_{t_0,t_0}-x_\infty)\cosh\!\big(a^2(t-t_0)\big)
%-
%x_{t_0,t_0}\sinh\!\big(a^2(t-t_0)\big)
%\Big].
%}
%\tag{B.3.4}
%\end{equation}
% ============================================================
\section{Relation to Conventional Interest Rates}
\label{appendix:C}

This appendix maps the instantaneous rolling forward rate (IRFR)
to the conventional rate $r_{t,T}$ with remaining maturity $T-t$.
Although the conventional rate fixes the maturity date $T$,
the IRFRs appearing in its definition continue to "roll'' with time.
We provide both discrete-time and continuous-time expressions and 
derive several key results used in Section~\ref{sec:econ}.

% ------------------------------------------------------------
\subsection{Expected Value and Instantaneous Variance of the Conventional Rate}
\label{appendix:C1}

In this section, key calculations and representations are mapped from the IRFR
onto the conventional interest rate \(r_{t,T}\) at time \(t\) with maturity \(T-t\).
The reader should note that the forward rates appearing in this representation remain rolling in maturity.

The IRFR relates to the conventional rate as
\[
r_{t,T}
=
\frac{1}{T-t}\sum_{S=t}^{T} x_{t,S}
\]
in discrete time, and
\begin{equation}
r_{t,T}
=
\frac{1}{T-t}\int_t^T x_{t,S}\,dS.
\tag{C.1.1}\label{eq:C11}
\end{equation}

\medskip
\noindent\textbf{Expected value.}

To compute the expected value of \(r_{t,T}\), we begin from
\begin{equation}
\mathbb{E}[x_{t,T}]
=
x_{t_0,t_0}e^{-a^2(T+t-2t_0)}
+
x_\infty\left(1-e^{-a^2(T-t_0)}\right).
\tag{C.1.2}\label{eq:C12}
\end{equation}

Integrating over maturities,
\[
\mathbb{E}\!\left[\int_t^T x_{t,S}\,dS\right]
=
\int_t^T
\left(
x_{t_0,t_0}e^{-a^2(S+t-2t_0)}
+
x_\infty\left(1-e^{-a^2(S-t_0)}\right)
\right)dS.
\]

Evaluating term by term,
\begin{align*}
\int_t^T x_{t_0,t_0}e^{-a^2(S+t-2t_0)}\,dS
&=
\frac{x_{t_0,t_0}}{a^2}
\left(
e^{-2a^2(t-t_0)}
-
e^{-a^2(T+t-2t_0)}
\right), \\
\int_t^T x_\infty\,dS
&=
x_\infty (T-t), \\
\int_t^T x_\infty e^{-a^2(S-t_0)}\,dS
&=
\frac{x_\infty}{a^2}
\left(
e^{-a^2(t-t_0)}
-
e^{-a^2(T-t_0)}
\right).
\end{align*}

Hence,
\begin{align}
\mathbb{E}\!\left[\int_t^T x_{t,S}\,dS\right]
&=
\frac{x_{t_0,t_0}}{a^2}
\left(
e^{-2a^2(t-t_0)}
-
e^{-a^2(T+t-2t_0)}
\right) \notag \\
&\quad
+ x_\infty (T-t)
- \frac{x_\infty}{a^2}
\left(
e^{-a^2(t-t_0)}
-
e^{-a^2(T-t_0)}
\right).
\tag{C.1.3}\label{eq:C13}
\end{align}

Using \eqref{eq:C11}, we obtain
\begin{equation}
\mathbb{E}[r_{t,T}]
=
x_\infty
+
\frac{1}{a^2(T-t)}
\left[
x_{t_0,t_0}
\left(
e^{-2a^2(t-t_0)}
-
e^{-a^2(T+t-2t_0)}
\right)
-
x_\infty
\left(
e^{-a^2(t-t_0)}
-
e^{-a^2(T-t_0)}
\right)
\right].
\tag{C.1.4}\label{eq:C14}
\end{equation}

\medskip
\noindent\textbf{Instantaneous variance.}

To compute the instantaneous variance, consider two maturities \(T_1\) and \(T_2\),
and define
\[
x_{t,T_1}
=
(x_{t,t}-x_\infty)e^{-a^2(T_1-t)} + x_\infty,
\qquad
y_{t,T_2}
=
(x_{t,t}-x_\infty)e^{-a^2(T_2-t)} + x_\infty.
\]

Then
\[
y_{t,T_2} - x_\infty
=
(x_{t,T_1}-x_\infty)e^{-a^2(T_2-T_1)}.
\]

Differentiating,
\[
\frac{\partial y_{t,T_2}}{\partial x_{t,T_1}}
=
e^{-a^2(T_2-T_1)}.
\]

Using covariance identities,
\[
\frac{\partial}{\partial t}
\mathrm{Cov}(x_{t,T_1},y_{t,T_2})
=
e^{-a^2(T_2-T_1)}
\frac{\partial}{\partial t}
\mathrm{Var}(x_{t,T_1}).
\]

Using the variance identity,
\[
\mathrm{Var}(x+y)
=
\mathrm{Var}(x)+\mathrm{Var}(y)+2\mathrm{Cov}(x,y),
\]
we obtain
\[
\frac{\partial}{\partial t}
\mathrm{Var}(x_{t,T_1}+y_{t,T_2})
=
\frac{\partial}{\partial t}
\mathrm{Var}(x_{t,T_1})
\left(1+e^{-a^2(T_2-T_1)}\right)^2.
\]

From the IRFR variance result,
\[
\frac{\partial}{\partial t}
\mathrm{Var}(x_{t,T_1})
=
2a^2
\left(
\left(x_{t_0,t_0}-\frac{x_\infty}{2}\right)e^{-a^2(t-t_0)}
+
\frac{x_\infty}{2}e^{a^2(t-t_0)}
\right)^2
e^{-2a^2(T_1-t_0)}.
\]

Hence, for the continuous case,
\begin{align}
\frac{\partial}{\partial t}
\mathrm{Var}\!\left(\int_t^T x_{t,S}\,dS\right)
&=
2a^2
\left(
\left(x_{t_0,t_0}-\frac{x_\infty}{2}\right)e^{-a^2(t-t_0)}
+
\frac{x_\infty}{2}e^{a^2(t-t_0)}
\right)^2 \notag \\
&\quad\times
\left(\int_t^T e^{-a^2(S-t_0)}\,dS\right)^2.
\tag{C.1.5}\label{eq:C15}
\end{align}

Using \eqref{eq:C11}, we obtain
\begin{equation}
\frac{\partial}{\partial t}\mathrm{Var}(r_{t,T})
=
\frac{2}{a^2(T-t)^2}
\left(
\left(x_{t_0,t_0}-\frac{x_\infty}{2}\right)e^{-a^2(t-t_0)}
+
\frac{x_\infty}{2}e^{a^2(t-t_0)}
\right)^2
\left(
e^{-a^2(t-t_0)}-e^{-a^2(T-t_0)}
\right)^2.
\tag{C.1.6}\label{eq:C16}
\end{equation}

An alternative representation is
\begin{equation}
\frac{\partial}{\partial t}\mathrm{Var}(r_{t,T})
=
\frac{2}{a^2(T-t)^2}
\left(
\left(x_{t_0,t_0}-\frac{x_\infty}{2}\right)e^{-2a^2(t-t_0)}
+\frac{x_\infty}{2}
\right)^2
\left(1-e^{-a^2(T-t)}\right)^2.
\tag{C.1.7}\label{eq:C17}
\end{equation}

% ------------------------------------------------------------

% ------------------------------------------------------------

%============----------------------------====================
% Appendix D
%=======================================================


\section{No-Arbitrage Justification of Rolling Maturity Invariance}
\label{Appedix D}

\subsection{Objective}
\label{D.1}

This appendix provides a rigorous no-arbitrage justification for Axiom 2 (Rolling Maturity Invariance). We show that if the stochastic dynamics of the instantaneous rolling forward rate (IRFR) depend on calendar time rather than solely on remaining maturity, then a self-financing trading strategy exists that generates arbitrage.

\subsection{IRFR as a Tradable P\&L Object}
\label{D.2}

Recall the bond pricing identity:
\[
\int_t^T x_{t,S} \, dS = -\ln P(t,T).
\]

Differentiating with respect to time implies that infinitesimal changes in the IRFR correspond to changes in zero-coupon bond prices. Therefore, there exists a self-financing portfolio of bonds whose instantaneous profit and loss satisfies:
\[
dV_t = dx_{t,T}.
\]

Thus, increments of the IRFR are directly associated with tradable P\&L.

\subsection{Two Equivalent Trading Strategies}
\label{D.3}

Fix a maturity horizon $\tau > 0$. Consider two strategies initiated at time $t$:

\paragraph{Strategy A (Hold Strategy)}
\begin{itemize}
\item Enter a forward exposure corresponding to $x_{t,T}$, where $T = t + \tau$
\item Hold continuously over $[t, t + \Delta t]$
\end{itemize}

The P\&L over the interval is:
\[
\Delta V^{\text{hold}} = x_{t+\Delta t, T} - x_{t,T}.
\]

\paragraph{Strategy B (Rolling Strategy)}
\begin{itemize}
\item Enter the same exposure at time $t$
\item At time $t + \Delta t$, close the position
\item Re-enter a forward exposure with the same remaining maturity $\tau$, i.e. $x_{t+\Delta t, T + \Delta t}$
\end{itemize}

The P\&L over the same interval is:
\[
\Delta V^{\text{roll}} = x_{t+\Delta t, T + \Delta t} - x_{t,T}.
\]

\subsection{No-Arbitrage Requirement}
\label{D.4}

Both strategies represent exposure to a forward rate with identical remaining maturity $\tau$. Therefore, no-arbitrage requires that they produce identical gains:
\[
\Delta V^{\text{hold}} = \Delta V^{\text{roll}}.
\]

Substituting expressions yields:
\[
x_{t+\Delta t, T} = x_{t+\Delta t, T + \Delta t}.
\]

Taking the limit as $\Delta t \to 0$, this implies:
\[
dx_{t,T} = dx_{t+dt, T+dt}
\]
in distribution.

\subsection{Implication for Dynamics}
\label{D.5}

The above identity shows that the stochastic evolution of the IRFR must be invariant along trajectories of constant remaining maturity:
\[
\tau = T - t.
\]

Therefore, the drift and diffusion coefficients must satisfy:
\[
\mu(x,t,T) = \mu(x, \tau), \quad \sigma(x,t,T) = \sigma(x, \tau).
\]

This establishes Axiom 2.

\subsection{Arbitrage Construction Under Violation}
\label{D.6)}

Suppose instead that dynamics depend explicitly on calendar time, so that:
\[
\mu(x,t,T) \neq \mu(x,t+\Delta t, T+\Delta t).
\]

Then:
\[
\mathbb{E}[\Delta V^{\text{hold}}] \neq \mathbb{E}[\Delta V^{\text{roll}}].
\]

Construct a portfolio:
\begin{itemize}
\item Long Strategy B (rolling)
\item Short Strategy A (hold)
\end{itemize}

This portfolio:
\begin{itemize}
\item has zero initial cost
\item eliminates common risk exposure
\item yields strictly positive expected return
\end{itemize}

Since both strategies are replicable via bond portfolios, this constitutes an arbitrage.

\subsection{Interpretation}
\label{D.7}

The result enforces a path-independence condition:
\begin{quote}
The P\&L from maintaining a forward exposure must be independent of whether the exposure is implemented continuously or via rolling.
\end{quote}

Equivalently, the stochastic law governing the IRFR must be invariant under the transformation:
\[
(t, T) \mapsto (t + \Delta t, T + \Delta t).
\]

This invariance is precisely the content of Axiom 2.

\subsection{Conclusion}
\label{D.8}

We have shown that rolling--maturity invariance is not merely a modelling convenience but a necessary condition for the absence of dynamic arbitrage. Any dependence of IRFR dynamics on calendar time would permit the construction of self-financing trading strategies that generate arbitrage through rolling inconsistencies.

Thus, Axiom 2 is justified as a no-arbitrage condition.

\end{appendices}

\end{document}
